% Lesson 43 — Reading: Texts about taking part in polls/surveys — Anglais Tle A — Cours signé OnBuch+
% Programme officiel MINESEC. Compiler avec : tectonic EN43-reading-texts-about-taking-part-in-polls-surveys.tex
\newcommand{\DOCMATIERE}{Anglais}
\newcommand{\DOCNIVEAU}{Tle A}
\newcommand{\DOCMODULE}{Chapter 4 : Citizenship and Human Rights}
\newcommand{\DOCLECON}{EN43}
\newcommand{\DOCTITRE}{Lesson 43 — Reading: Texts about taking part in polls/surveys}
% OnBuch+ — préambule « cours signés » ENRICHI (compilé avec tectonic / XeLaTeX)
% Système visuel « Pop » d'OnBuch+ : fond crème (jamais blanc), bordures encre
% épaisses, ombres « dures » décalées, titres Archivo Black en capitales.
% Version MATHÉMATIQUES : graphes (pgfplots/tikz), boîtes pédagogiques
% (définition, propriété, méthode, attention, le savais-tu, exercice résolu,
% exercice, TP, bilan), page de garde et signature OnBuch+ ancrée en bas.
%
% Macros attendues dans main.tex AVANT \input{preamble} :
%   \DOCMATIERE  \DOCNIVEAU  \DOCMODULE  \DOCLECON (numéro)  \DOCTITRE
\documentclass[11pt]{article}

\usepackage{fontspec}
\usepackage[english,french]{babel} % français = langue principale ; l'anglais via \eng{..} / textebox
\usepackage[a4paper,margin=2cm,top=3.4cm,bottom=2.8cm,headheight=1.4cm,footskip=1.3cm]{geometry}
\usepackage{amsmath,amssymb,amsfonts}
\usepackage{mathtools} % \xmapsto, \coloneqq... souvent utilisés par les modèles
\usepackage{cancel} % simplifications barrées (bilans de forces)
% \abs{z} (module, valeur absolue) et \norm{u} (norme), utilisés par les modèles
\providecommand{\abs}{}\let\abs\relax\DeclarePairedDelimiter{\abs}{\lvert}{\rvert}
\providecommand{\norm}{}\let\norm\relax\DeclarePairedDelimiter{\norm}{\lVert}{\rVert}
\usepackage[table]{xcolor} % [table] : fournit aussi \rowcolors (alternance de lignes)
\usepackage{pagecolor}
\usepackage{tikz}
\usetikzlibrary{arrows.meta,positioning,calc,shapes.geometric,shapes.misc,shapes.arrows,decorations.pathmorphing,decorations.pathreplacing,decorations.markings,patterns,shadows,mindmap,trees,fit,backgrounds,matrix,intersections,angles,quotes}
\usepackage{pgfplots}
\usepackage[version=4]{mhchem} % équations nucléaires \ce{^{235}_{92}U}
\usepackage[european,siunitx]{circuitikz} % schémas électriques
\pgfplotsset{compat=1.18}
\usepgfplotslibrary{fillbetween,groupplots} % aires entre courbes (calcul intégral)
\usepackage{enumitem}
\usepackage{fancyhdr}
% [most] charge toutes les bibliothèques tcolorbox (listings, minted, raster,
% documentation...) et gonfle inutilement la mémoire TeX sur un document long
% (~300 boîtes) : on ne charge que ce qui est réellement utilisé.
\usepackage[skins,breakable]{tcolorbox}
\usepackage{graphicx}
\usepackage{colortbl}
\usepackage{booktabs}
\usepackage{tabularx}
\usepackage{diagbox} % case d'en-tête diagonale des tableaux à double entrée
% Colonnes extensibles centrée / gauche / droite (C, L, R), souvent utilisées par les modèles
\newcolumntype{C}{>{\centering\arraybackslash}X}
\newcolumntype{L}{>{\raggedright\arraybackslash}X}
\newcolumntype{R}{>{\raggedleft\arraybackslash}X}
\usepackage{array}
\usepackage{multicol}
\usepackage{multirow}
\usepackage{siunitx}
% parse-numbers=false : accepte aussi \SI{2,5 \times 10^{-3}}{...}
\sisetup{locale=FR,output-decimal-marker={,},group-minimum-digits=4,parse-numbers=false}
\DeclareSIUnit\ohm{\ensuremath{\symup{\Omega}}} % U+2126 absent de la police
\DeclareSIUnit\division{div} % réglages d'oscilloscope (ms/div, V/div)
\DeclareSIUnit\tr{tr} % tours (vitesse de rotation, ex. \SI{1500}{\tr\per\minute})
\DeclareSIUnit\molar{M} % concentration molaire (ex. \SI{3}{\molar}, \SI{1}{\micro\molar})
\DeclareSIUnit\an{an} % année (le modèle confond parfois avec \year, un compteur TeX) — ex. \SI{5}{\kilo\meter\per\an}
\DeclareSIUnit\FCFA{FCFA} % franc CFA (ex. \SI{500}{\FCFA\per\kilo\gram})
\DeclareSIUnit\degreeBrix{\textdegree Brix} % teneur en sucre d'un sirop/jus (réfractométrie)
\DeclareSIUnit\month{mois} % le modèle utilise parfois \month, un compteur TeX, comme unité de durée
\usepackage{qrcode}
% \degree : commande fréquemment utilisée par les modèles pour un angle ou une
% température en dehors de \SI{}{} (non fournie par les paquets chargés).
\providecommand{\degree}{^{\circ}}
% \qed (fin de démonstration), utilisé par les modèles sans amsthm
\providecommand{\qed}{\hfill$\square$}
% \barre : double barre des valeurs interdites dans un tableau de variations (tkz-tab)
\providecommand{\barre}{\ensuremath{\|}}
% environnement proof (amsthm non chargé) utilisé par les modèles
\ifdefined\proof\else\newenvironment{proof}[1][Démonstration]{\par\noindent\textit{#1.}\ }{\qed\par}\fi
\usepackage{lastpage}
\usepackage{hyperref}
\hypersetup{hidelinks}

% ============================================================
% Palette « Pop » OnBuch+ — mêmes hex que la classe P (plus_ui.dart)
% ============================================================
\definecolor{poporange}{HTML}{F59321}
\definecolor{popdark}{HTML}{E07A0C}
\definecolor{popcream}{HTML}{FFF6E8}
\definecolor{popink}{HTML}{1C1714}
\definecolor{poppurple}{HTML}{7A5AE0}
\definecolor{popgreen}{HTML}{2E9E6A}
\definecolor{poppink}{HTML}{E0457B}
\definecolor{popblue}{HTML}{2F7DE1}
\definecolor{popgold}{HTML}{C98A12}
\definecolor{popmuted}{HTML}{8A7F76}
\definecolor{popline}{HTML}{EADFD2}
\definecolor{popgray}{HTML}{8A7F76}
\colorlet{popgrey}{popgray}
% Alias tolérants (le modèle nomme parfois les couleurs différemment)
\colorlet{popwhite}{white}
\colorlet{popblack}{popink}
\colorlet{popred}{poppink}
\colorlet{popyellow}{popgold}
% Teintes claires (fonds de tableaux/encadrés) — le modèle nomme parfois
% les couleurs avec un suffixe L (ex. popblueL) sans les avoir définies.
\colorlet{poporangeL}{poporange!20}
\colorlet{popdarkL}{popdark!20}
\colorlet{poppurpleL}{poppurple!20}
\colorlet{popgreenL}{popgreen!20}
\colorlet{poppinkL}{poppink!20}
\colorlet{popblueL}{popblue!20}
\colorlet{popgoldL}{popgold!20}
\colorlet{popredL}{popred!20}
\colorlet{popgrayL}{popgray!20}
% Teintes claires pour fonds de tableaux / zones
\colorlet{popblueL}{popblue!10!white}
\colorlet{poporangeL}{poporange!14!white}
\colorlet{popgreenL}{popgreen!12!white}
\colorlet{poppurpleL}{poppurple!12!white}
\colorlet{poppinkL}{poppink!12!white}
\colorlet{popgoldL}{popgold!14!white}

\pagecolor{popcream}

% ============================================================
% Typographie
% ============================================================
\defaultfontfeatures{Ligatures=TeX}
\setmainfont{PlusJakartaSans-Medium.ttf}[Path=../../../fonts/,BoldFont=PlusJakartaSans-Bold.ttf]
% Symboles, API et emojis absents de la police principale : repli sur DejaVu Sans (caractères actifs)
\newfontface\symfont{DejaVuSans.ttf}[Path=../../../fonts/]
\makeatletter
\newcommand\symactive[1]{\catcode#1=13 \begingroup\lccode`\~=#1 \lowercase{\endgroup\def~}{{\symfont\char#1}}}
\newcommand\symblank[1]{\catcode#1=13 \begingroup\lccode`\~=#1 \lowercase{\endgroup\def~}{}}
\newcommand\symhyph[1]{\catcode#1=13 \begingroup\lccode`\~=#1 \lowercase{\endgroup\def~}{\mbox{-}}}
\newcount\symc
\newcommand\symrange[2]{\symc=#1 \@whilenum\symc<#2 \do{\expandafter\symactive\expandafter{\the\symc}\advance\symc by 1 }}
\newcommand\symblankrange[2]{\symc=#1 \@whilenum\symc<#2 \do{\expandafter\symblank\expandafter{\the\symc}\advance\symc by 1 }}
\makeatother
\symrange{"0250}{"02FF}\symactive{"2713}\symactive{"2714}\symactive{"2717}\symactive{"2718}\symactive{"2610}\symactive{"2611}\symactive{"2612}
\symhyph{"2011}\symhyph{"2E1A}\symblankrange{"1F300}{"1FAFF}\symblank{"FE0F}
% Maths en sans-serif (Fira Math) avec chiffres et lettres droites de Plus
% Jakarta, pour que les formules se fondent dans le texte.
\usepackage[math-style=ISO,bold-style=ISO]{unicode-math}
\setmathfont{FiraMath-Regular.otf}[Path=../../../fonts/]
\setmathfont{PlusJakartaSans-Medium.ttf}[Path=../../../fonts/,range={up/num,up/{Latin,latin},\mathup}]
\usepackage{icomma} % virgule décimale sans espace en mode math
% Caractères mathématiques Unicode absents des polices de texte (Plus Jakarta,
% Archivo Black) : rendus par leur équivalent mathématique, en texte comme en
% formule — sinon ils s'affichent en carré vide (ex. « Arithmétique dans ℤ »).
\usepackage{newunicodechar}
\newunicodechar{ℝ}{\ensuremath{\mathbb{R}}}
\newunicodechar{ℤ}{\ensuremath{\mathbb{Z}}}
\newunicodechar{ℂ}{\ensuremath{\mathbb{C}}}
\newunicodechar{ℕ}{\ensuremath{\mathbb{N}}}
\newunicodechar{ℚ}{\ensuremath{\mathbb{Q}}}
\newunicodechar{𝔻}{\ensuremath{\mathbb{D}}}
\newunicodechar{α}{\ensuremath{\alpha}}
\newunicodechar{β}{\ensuremath{\beta}}
\newunicodechar{γ}{\ensuremath{\gamma}}
\newunicodechar{δ}{\ensuremath{\delta}}
\newunicodechar{ε}{\ensuremath{\varepsilon}}
\newunicodechar{θ}{\ensuremath{\theta}}
\newunicodechar{λ}{\ensuremath{\lambda}}
\newunicodechar{μ}{\ensuremath{\mu}}
\newunicodechar{σ}{\ensuremath{\sigma}}
\newunicodechar{τ}{\ensuremath{\tau}}
\newunicodechar{φ}{\ensuremath{\varphi}}
\newunicodechar{ψ}{\ensuremath{\psi}}
\newunicodechar{ω}{\ensuremath{\omega}}
\newunicodechar{Σ}{\ensuremath{\Sigma}}
% majuscules produites par \MakeUppercase dans les titres (α-aminés -> Α)
\newunicodechar{Α}{\ensuremath{\alpha}}
\newunicodechar{Β}{\ensuremath{\beta}}
\newunicodechar{Γ}{\ensuremath{\Gamma}}
\newunicodechar{Δ}{\ensuremath{\Delta}}
\newunicodechar{Ε}{\ensuremath{\varepsilon}}
\newunicodechar{Θ}{\ensuremath{\Theta}}
\newunicodechar{Λ}{\ensuremath{\Lambda}}
\newunicodechar{Μ}{\ensuremath{\mu}}
\newunicodechar{Τ}{\ensuremath{\tau}}
\newunicodechar{Φ}{\ensuremath{\Phi}}
\newunicodechar{Ψ}{\ensuremath{\Psi}}
\newunicodechar{Ω}{\ensuremath{\Omega}}
\newunicodechar{↦}{\ensuremath{\mapsto}}
\newunicodechar{∈}{\ensuremath{\in}}
\newunicodechar{∉}{\ensuremath{\notin}}
\newunicodechar{∧}{\ensuremath{\wedge}}
\newunicodechar{⇔}{\ensuremath{\Leftrightarrow}}
\newunicodechar{⟺}{\ensuremath{\Longleftrightarrow}}
\newunicodechar{⇒}{\ensuremath{\Rightarrow}}
\newunicodechar{⟹}{\ensuremath{\Longrightarrow}}
\newunicodechar{∘}{\ensuremath{\circ}}
\newunicodechar{∪}{\ensuremath{\cup}}
\newunicodechar{∩}{\ensuremath{\cap}}
\newunicodechar{≡}{\ensuremath{\equiv}}
\newunicodechar{⊂}{\ensuremath{\subset}}
\newunicodechar{⊥}{\ensuremath{\perp}}
\newunicodechar{∃}{\ensuremath{\exists}}
\newunicodechar{∀}{\ensuremath{\forall}}
\newunicodechar{∛}{\ensuremath{\sqrt[3]{\,}}}
\newunicodechar{∜}{\ensuremath{\sqrt[4]{\,}}}
\newunicodechar{⃗}{\ensuremath{{}^{\to}}}
\newunicodechar{ⁿ}{\ifmmode^{n}\else\textsuperscript{\ensuremath{n}}\fi}
\newunicodechar{ˣ}{\ifmmode^{x}\else\textsuperscript{\ensuremath{x}}\fi}
\newunicodechar{ᵇ}{\ifmmode^{b}\else\textsuperscript{\ensuremath{b}}\fi}
\newunicodechar{ʳ}{\ifmmode^{r}\else\textsuperscript{\ensuremath{r}}\fi}
\newunicodechar{ᵘ}{\ifmmode^{u}\else\textsuperscript{\ensuremath{u}}\fi}
\newunicodechar{ʸ}{\ifmmode^{y}\else\textsuperscript{\ensuremath{y}}\fi}
\newunicodechar{ᵃ}{\ifmmode^{a}\else\textsuperscript{\ensuremath{a}}\fi}
\newunicodechar{ᵉ}{\ifmmode^{e}\else\textsuperscript{\ensuremath{e}}\fi}
\newunicodechar{ᵏ}{\ifmmode^{k}\else\textsuperscript{\ensuremath{k}}\fi}
\newunicodechar{ᵖ}{\ifmmode^{p}\else\textsuperscript{\ensuremath{p}}\fi}
\newunicodechar{ᴺ}{\ifmmode^{N}\else\textsuperscript{\ensuremath{N}}\fi}
\newunicodechar{ᶜ}{\ifmmode^{c}\else\textsuperscript{\ensuremath{c}}\fi}
\newunicodechar{ʰ}{\ifmmode^{h}\else\textsuperscript{\ensuremath{h}}\fi}
\newunicodechar{ᵠ}{\ifmmode^{q}\else\textsuperscript{\ensuremath{q}}\fi}
\newunicodechar{ₙ}{\ifmmode_{n}\else\textsubscript{\ensuremath{n}}\fi}
\newunicodechar{ₐ}{\ifmmode_{a}\else\textsubscript{\ensuremath{a}}\fi}
\newunicodechar{ᵢ}{\ifmmode_{i}\else\textsubscript{\ensuremath{i}}\fi}
\newunicodechar{ᵣ}{\ifmmode_{r}\else\textsubscript{\ensuremath{r}}\fi}
\newunicodechar{ₖ}{\ifmmode_{k}\else\textsubscript{\ensuremath{k}}\fi}
\newunicodechar{ᵦ}{\ifmmode_{\beta}\else\textsubscript{\ensuremath{\beta}}\fi}
\newunicodechar{ₓ}{\ifmmode_{x}\else\textsubscript{\ensuremath{x}}\fi}
\newfontfamily\popdisplay{ArchivoBlack-Regular.ttf}[Path=../../../fonts/]
\newfontfamily\poplogo{SpaceGrotesk-Bold.ttf}[Path=../../../fonts/]
\newfontfamily\popsemi{PlusJakartaSans-SemiBold.ttf}[Path=../../../fonts/]
\newfontfamily\popxbold{PlusJakartaSans-ExtraBold.ttf}[Path=../../../fonts/]
\newfontfamily\popxboldls{PlusJakartaSans-ExtraBold.ttf}[Path=../../../fonts/,LetterSpace=3.0]

\setlist[enumerate]{leftmargin=1.7em,itemsep=5pt,topsep=4pt}
\setlist[itemize]{leftmargin=1.7em,itemsep=4pt,topsep=4pt,label={\color{poporange}\textbullet}}
\setlength{\parindent}{0pt}
\setlength{\parskip}{5pt}
\renewcommand{\arraystretch}{1.25}

% pgfplots : style Pop par défaut
\pgfplotsset{
  every axis/.append style={axis line style={popink,line width=1pt},tick style={popink},
    grid style={popline,line width=0.5pt},label style={font=\small},tick label style={font=\footnotesize}},
  popcurve/.style={poporange,line width=1.8pt,no markers,smooth},
}
% Même style disponible dans un \draw TikZ simple (hors pgfplots)
\tikzset{popcurve/.style={poporange,line width=1.8pt}}
\tikzset{popgrid/.style={popline,very thin}}
\colorlet{popcurve}{poporange} % le nom du style est parfois utilisé comme couleur

% ============================================================
% Pill / squiggle
% ============================================================
\newcommand{\poppill}[3]{%
  \tikz[baseline=-0.55ex]\node[draw=popink,line width=1.2pt,rounded corners=2.6pt,%
    fill=#2,inner xsep=6.5pt,inner ysep=3pt]{%
    \popxboldls\fontsize{7}{7}\selectfont\color{#3}\MakeUppercase{#1}};%
}
\newcommand{\popsquiggle}[1]{%
  \tikz[baseline=-0.4ex]{
    \draw[#1,line width=2.6pt,line cap=round]
      plot [smooth] coordinates {(0,0.09) (0.34,0.01) (0.68,0.21) (1.02,0.01) (1.36,0.10)};
  }%
}
% Mot-clé surligné (marqueur orange sous le texte)
\newcommand{\cle}[1]{{\popxbold\color{popdark}#1}}

% ============================================================
% En-tête / pied
% ============================================================
\pagestyle{fancy}
\fancyhf{}
\renewcommand{\headrulewidth}{0pt}
\renewcommand{\footrulewidth}{0pt}
\lhead{\raisebox{-0.15ex}{\poplogo\fontsize{13}{13}\selectfont\color{popink}OnBuch\color{poporange}+}}
\rhead{\poppill{\DOCMATIERE\ \textbullet\ \DOCNIVEAU\ \textbullet\ Leçon \DOCLECON}{white}{popink}}
\lfoot{\popxbold\fontsize{7.6}{7.6}\selectfont\color{popmuted}\MakeUppercase{Cours signé OnBuch+}\ \textbullet\ {\popsemi\DOCTITRE}}
\rfoot{\tikz[baseline=-0.6ex]\node[rounded corners=8.5pt,fill=popink,minimum height=17pt,minimum width=17pt,inner xsep=5pt,inner sep=0]{\popxbold\fontsize{8}{8}\selectfont\color{popcream}\thepage\,/\,\pageref*{LastPage}};}

% ============================================================
% Titres
% ============================================================
\newcommand{\chaptitre}[2]{%
  \begin{tcolorbox}[enhanced,colback=poporange,colframe=popink,boxrule=2.6pt,arc=11pt,
      drop shadow=popink,left=15pt,right=15pt,top=12pt,bottom=11pt,before skip=3pt,after skip=18pt]
    \poppill{#2}{popink}{popcream}\par\vspace{9pt}
    {\raggedright\hyphenpenalty=10000\exhyphenpenalty=10000\popdisplay\fontsize{22}{24}\selectfont\color{popcream}\MakeUppercase{#1}\par}
  \end{tcolorbox}
}
% Section numérotée : pastille orange avec le numéro + titre Archivo
\newcounter{popsec}
\newcommand{\coursec}[1]{%
  \stepcounter{popsec}\setcounter{popsub}{0}%
  \par\vspace{16pt}\noindent
  \begin{minipage}{\linewidth}
  \tikz[baseline=-0.7ex]\node[draw=popink,line width=1.6pt,fill=poporange,rounded corners=4pt,minimum size=22pt,inner sep=0]{\popdisplay\fontsize{12}{12}\selectfont\color{popcream}\thepopsec};%
  \hspace{8pt}{\popdisplay\fontsize{15}{17}\selectfont\color{popink}\MakeUppercase{#1}}\par
  \vspace{4pt}\noindent{\color{poporange}\rule{36pt}{3.6pt}}
  \end{minipage}\par\nopagebreak\vspace{8pt}
}
\newcounter{popsub}
\newcommand{\courssub}[1]{%
  \stepcounter{popsub}%
  \par\vspace{9pt}\noindent{\popxbold\fontsize{11.5}{13}\selectfont\color{popink}{\color{poporange}\thepopsec.\thepopsub}\hspace{6pt}#1}\par\nopagebreak\vspace{4pt}
}

% ============================================================
% Boîtes pédagogiques — même recette : fond blanc, bordure épaisse colorée,
% ombre dure encre, badge plein. Titre optionnel [#1].
% ============================================================
\tcbset{popbase/.style={breakable,enhanced,colback=white,boxrule=2.3pt,arc=9pt,
  shadow={3pt}{-3pt}{0pt}{popink},left=14pt,right=14pt,top=11pt,bottom=11pt,before skip=11pt,after skip=15pt,
  fontupper=\popsemi\fontsize{10.6}{14.5}\selectfont\color{popink},
  fontlower=\popsemi\fontsize{10.6}{14.5}\selectfont\color{popink}}}
% Coche dessinée (le glyphe ✓ n'existe pas dans les polices du document)
\newcommand{\popcheck}[1]{\tikz[baseline=-0.5ex]\draw[#1,line width=1.6pt,line cap=round,line join=round](-0.13,0.0)--(-0.04,-0.09)--(0.13,0.1);}
\providecommand{\coche}{\popcheck{popgreen}}
\providecommand{\resultat}[1]{\par\smallskip\noindent\fcolorbox{popgreen}{popgreenL}{\parbox{\dimexpr\linewidth-2\fboxsep-2\fboxrule\relax}{\textbf{Résultat.} #1}}\par\smallskip}
\newcommand{\popboxhead}[3]{\poppill{#1}{#2}{white}\ifx\relax#3\relax\else\hspace{6pt}{\popxbold\fontsize{10.6}{12}\selectfont\color{popink}#3}\fi\par\vspace{7pt}}

\newtcolorbox{aretenir}[1][]{popbase,colframe=popgreen,before upper={\popboxhead{À retenir}{popgreen}{#1}}}
% Texte en anglais (document de lecture, dialogue, extrait) : cadre
% d'étude. Un titre optionnel (source, auteur) s'affiche dans l'en-tête.
\newtcolorbox{textebox}[1][]{popbase,colframe=popblue,before upper={\popboxhead{Text}{popblue}{#1}\begin{otherlanguage}{english}},after upper={\end{otherlanguage}}}
% Marques ✓ / ✗ (forme correcte / fautive) souvent utilisées par les modèles
\providecommand{\cross}{\textcolor{poppink}{\ensuremath{\times}}}
\providecommand{\xmark}{\cross}\providecommand{\crossmark}{\cross}\providecommand{\cmark}{\ensuremath{\checkmark}}
% Ligne de réponse / justification à compléter (inventée par les modèles)
\providecommand{\justification}[1]{\par\noindent\textit{Justification~:} \dotfill\par}
% \textipa (package tipa non chargé) : la police ne couvre pas l'API -> simple passage
\providecommand{\textipa}[1]{#1}
% Soulignement (ulem non chargé) : \uline, \ul
\providecommand{\uline}[1]{\underline{#1}}\providecommand{\ul}[1]{\underline{#1}}
% \glossaire{\item ...} inventée par les modèles : liste de glossaire
\providecommand{\glossaire}[1]{\begin{itemize}#1\end{itemize}}
% \alert (beamer) inventée par les modèles : texte mis en évidence
\providecommand{\alert}[1]{\textbf{\textcolor{poppink}{#1}}}
% variantes ASCII d'ordinaux écrites en macro
\providecommand{\iere}{\textsuperscript{re}}\providecommand{\ieme}{\textsuperscript{e}}\providecommand{\ere}{\textsuperscript{re}}\providecommand{\eme}{\textsuperscript{e}}\providecommand{\er}{\textsuperscript{er}}\providecommand{\ie}{\textsuperscript{e}}
% Ordinaux français écrits en macro par les modèles : 1\ière, 3\ième
\newcommand{\ière}{\textsuperscript{re}}\newcommand{\ième}{\textsuperscript{e}}
% \exemple (inventée par les modèles) : étiquette « Exemple : »
\providecommand{\exemple}{\textbf{Exemple~:}\ }
% \square utilisable aussi hors mode math (cases à cocher)
\AtBeginDocument{\let\squareMath\square\renewcommand{\square}{\ensuremath{\squareMath}}}
% Mot ou phrase en anglais à l'intérieur d'un paragraphe français
\newcommand{\eng}[1]{\ifmmode\text{\textit{#1}}\else\begingroup\selectlanguage{english}\itshape #1\endgroup\fi}
\let\esp\eng % alias toléré
% flèches d'intonation inventées par les modèles
\providecommand{\textsearrow}{}\renewcommand{\textsearrow}{\ensuremath{\searrow}}\providecommand{\textnearrow}{}\renewcommand{\textnearrow}{\ensuremath{\nearrow}}
% \fr{..} : mot français (inventé par les modèles)
\providecommand{\fr}[1]{\textit{#1}}
\newtcolorbox{exemplebox}[1][]{popbase,colframe=poporange,before upper={\popboxhead{Exemple}{poporange}{#1}}}
\newtcolorbox{definition}[1][]{popbase,colframe=popblue,before upper={\popboxhead{Définition}{popblue}{#1}}}
\newtcolorbox{propriete}[1][]{popbase,colframe=poppurple,before upper={\popboxhead{Propriété}{poppurple}{#1}}}
\newtcolorbox{methode}[1][]{popbase,colframe=popgold,before upper={\popboxhead{Méthode}{popgold}{#1}}}
\newtcolorbox{attention}[1][]{popbase,colframe=poppink,before upper={\popboxhead{Attention}{poppink}{#1}}}
\newtcolorbox{experience}[1][]{popbase,colframe=popblue,colback=popblueL,before upper={\popboxhead{Expérience}{popblue}{#1}}}
% « Le savais-tu ? » : carte violette pleine (contexte, culture, Cameroun)
\newtcolorbox{savaistu}[1][]{popbase,colback=poppurple,colframe=popink,
  fontupper=\popsemi\fontsize{10.4}{14.2}\selectfont\color{white},
  before upper={\poppill{Le savais-tu\,?}{white}{poppurple}\ifx\relax#1\relax\else\hspace{6pt}{\popxbold\color{white}#1}\fi\par\vspace{7pt}}}
% Exercice résolu : énoncé en haut, \tcblower puis la solution sur fond crème
\newcounter{popexo}\newcounter{popexr}
\newtcolorbox{exoresolu}[1][]{popbase,colframe=poporange,colbacklower=popcream,
  segmentation style={popink,line width=1.2pt,dashed},
  before upper={\stepcounter{popexr}\popboxhead{Exercice résolu \thepopexr}{poporange}{#1}},
  before lower={\poppill{Solution}{popink}{popcream}\par\vspace{6pt}}}
% Exercice (non résolu en place) — la correction est dans la partie Corrigés
\newtcolorbox{exercice}[1][]{popbase,colframe=popink,
  before upper={\stepcounter{popexo}\popboxhead{Exercice \thepopexo}{popink}{#1}}}
% Corrigé
\newtcolorbox{corrige}[1][]{popbase,colframe=popgreen,colback=popgreenL,
  before upper={\popboxhead{Corrigé}{popgreen}{#1}}}
% Objectifs / prérequis / bilan
\newtcolorbox{objectifs}{popbase,colframe=popink,colback=poporangeL,
  before upper={\popboxhead{Objectifs de la leçon}{popink}{}}}
\newtcolorbox{prerequis}{popbase,colframe=popink,colback=popblueL,
  before upper={\popboxhead{Prérequis}{popblue}{}}}
\newtcolorbox{bilan}{popbase,colframe=popink,colback=white,boxrule=2.8pt,
  before upper={\popboxhead{Fiche bilan}{poporange}{L'essentiel à connaître pour l'examen}}}
% Figure encadrée avec légende
\newtcolorbox{popfigure}[1][]{enhanced,breakable=false,colback=white,colframe=popline,boxrule=1.4pt,arc=8pt,
  left=10pt,right=10pt,top=10pt,bottom=8pt,before skip=10pt,after skip=12pt,halign=center,
  lower separated=false,halign lower=center,
  fontlower=\popsemi\fontsize{9}{11.5}\selectfont\color{popmuted},
  #1}
\newcommand{\legende}[1]{\par\vspace{6pt}{\popsemi\fontsize{9}{11.5}\selectfont\color{popmuted}#1}}

% ============================================================
% Signature OnBuch+
% ============================================================
\newcommand{\poplogoinline}[1]{{\poplogo\fontsize{#1}{#1}\selectfont\color{popink}OnBuch\color{poporange}+}}
\newcommand{\signatureonbuch}{%
  \begin{tcolorbox}[enhanced,colback=popink,colframe=popink,boxrule=2.2pt,arc=10pt,
      drop shadow=poporange,left=14pt,right=14pt,top=10pt,bottom=10pt,before skip=0pt,after skip=0pt]
    \tikz[baseline=-0.7ex]\node[circle,fill=popgreen,draw=popcream,line width=1.3pt,minimum size=22pt,inner sep=0]{\popcheck{white}};%
    \hspace{9pt}\begin{minipage}[c]{0.62\linewidth}
      {\popxbold\fontsize{12}{13}\selectfont\color{popcream}Signé\ \textbullet\ {\poplogo OnBuch\color{poporange}+}}\par\vspace{2pt}
      {\raggedright\popsemi\fontsize{8}{10.5}\selectfont\color{popline}\MakeUppercase{Équipe pédagogique OnBuch+}\\\MakeUppercase{Conforme au programme officiel MINESEC}\par}
    \end{minipage}\hfill
    {\poplogo\fontsize{22}{22}\selectfont\color{popcream}OnBuch\color{poporange}+}
  \end{tcolorbox}%
}

% ============================================================
% Page de garde
% #1 titre · #2 matière · #3 niveau · #4 module · #5 n° leçon · #6 durée ·
% #7 description / objectifs (texte ou itemize)
% ============================================================
\newcommand{\pagegarde}[7]{%
  \thispagestyle{empty}
  \newgeometry{a4paper,margin=1.7cm,top=1.6cm,bottom=1.4cm}%
  \begin{tikzpicture}[remember picture,overlay]
    \fill[poporange!18] ([xshift=-3.2cm,yshift=-3.6cm]current page.north east) circle (4.6cm);
    \fill[poppurple!14] ([xshift=2.4cm,yshift=7.6cm]current page.south west) circle (3.8cm);
  \end{tikzpicture}%
  {\poplogo\fontsize{30}{30}\selectfont\color{popink}OnBuch\color{poporange}+}\hfill
  \raisebox{0.6ex}{\poppill{Cours signé \textbullet\ Premium}{popink}{popcream}}\par
  \vspace{1pt}\popsquiggle{poppurple}\par
  \vspace{8pt}
  \begin{tcolorbox}[enhanced,colback=poporange,colframe=popink,boxrule=2.8pt,arc=13pt,
      drop shadow=popink,left=18pt,right=18pt,top=12pt,bottom=13pt,after skip=10pt]
    \poppill{#2 \textbullet\ #3}{popink}{popcream}\hspace{5pt}\poppill{#4}{white}{popink}\par\vspace{8pt}
    {\popdisplay\fontsize{14}{14}\selectfont\color{popink}LEÇON #5}\par\vspace{4pt}
    {\raggedright\hyphenpenalty=10000\exhyphenpenalty=10000\popdisplay\fontsize{25}{27}\selectfont\color{popcream}\MakeUppercase{#1}\par}
    \vspace{8pt}
    {\popxbold\fontsize{9.5}{9.5}\selectfont\color{popink}\MakeUppercase{Durée conseillée : #6}}
  \end{tcolorbox}
  \begin{tcolorbox}[enhanced,colback=white,colframe=popink,boxrule=2.2pt,arc=10pt,
      drop shadow=popink,left=16pt,right=16pt,top=10pt,bottom=10pt,after skip=8pt]
    \poppill{Dans cette leçon}{poppurple}{white}\par\vspace{5pt}
    \popsemi\fontsize{9.3}{12.6}\selectfont\color{popink}#7
  \end{tcolorbox}
  \noindent\begin{minipage}[t]{0.487\linewidth}
    \begin{tcolorbox}[enhanced,colback=popgreen,colframe=popink,boxrule=2.2pt,arc=10pt,
        drop shadow=popink,left=11pt,right=11pt,top=10pt,bottom=10pt,height=3.1cm,valign=center]
      \tcbox[colback=white,colframe=popink,boxrule=1.3pt,arc=5pt,boxsep=0pt,nobeforeafter,
        left=3pt,right=3pt,top=3pt,bottom=3pt,valign=center]{\qrcode[height=1.9cm]{https://play.google.com/store/apps/details?id=com.luvvix.onbuch&hl=fr}}%
      \hspace{8pt}\begin{minipage}[c]{0.5\linewidth}
        {\popdisplay\fontsize{11}{12.5}\selectfont\color{white}\MakeUppercase{Télécharge OnBuch}}\par\vspace{4pt}
        {\raggedright\popsemi\fontsize{8.4}{11}\selectfont\color{white}Gratuit \textbullet\ Google Play\\Tuteur IA, annales, concours, résultats\par}
      \end{minipage}
    \end{tcolorbox}
  \end{minipage}\hfill
  \begin{minipage}[t]{0.487\linewidth}
    \begin{tcolorbox}[enhanced,colback=poppurple,colframe=popink,boxrule=2.2pt,arc=10pt,
        drop shadow=popink,left=11pt,right=11pt,top=10pt,bottom=10pt,height=3.1cm,valign=center]
      \tikz[baseline=-0.7ex]\node[circle,fill=white,draw=popink,line width=1.3pt,minimum size=1.6cm,inner sep=0]{\popdisplay\fontsize{24}{24}\selectfont\color{poppurple}+};%
      \hspace{8pt}\begin{minipage}[c]{0.6\linewidth}
        {\popdisplay\fontsize{11}{12.5}\selectfont\color{white}\MakeUppercase{Passe à OnBuch+}}\par\vspace{4pt}
        {\raggedright\popsemi\fontsize{8.4}{11}\selectfont\color{white}Tous les cours signés, Tuteur IA illimité, fiches PDF — onglet OnBuch+ de l'app\par}
      \end{minipage}
    \end{tcolorbox}
  \end{minipage}
  \vspace{8pt}
  \popcontenu
  \vfill
  \signatureonbuch
  \restoregeometry
  \newpage
}

% Rangée « Ce cours contient » (page de garde)
\newcommand{\poptile}[3]{% #1 couleur #2 gros texte #3 légende
  \begin{tcolorbox}[enhanced,colback=#1,colframe=popink,boxrule=2pt,arc=9pt,shadow={2.5pt}{-2.5pt}{0pt}{popink},
      left=6pt,right=6pt,top=9pt,bottom=9pt,halign=center,valign=center,height=1.75cm,width=\linewidth,nobeforeafter]
    {\popdisplay\fontsize{12.5}{13}\selectfont\color{white}\MakeUppercase{#2}}\par\vspace{3pt}
    {\centering\popsemi\fontsize{7.8}{9.5}\selectfont\color{white}#3\par}
  \end{tcolorbox}}
\newcommand{\popcontenu}{%
  \par\noindent{\popxboldls\fontsize{7.5}{7.5}\selectfont\color{popmuted}CE COURS SIGNÉ CONTIENT}\par\vspace{7pt}
  \noindent\begin{minipage}[t]{0.235\linewidth}\poptile{popblue}{Cours}{complet, illustré, pas à pas}\end{minipage}\hfill
  \begin{minipage}[t]{0.235\linewidth}\poptile{poporange}{Méthodes}{et exercices résolus}\end{minipage}\hfill
  \begin{minipage}[t]{0.235\linewidth}\poptile{poppink}{Exercices}{type examen + corrigés}\end{minipage}\hfill
  \begin{minipage}[t]{0.235\linewidth}\poptile{popgreen}{Fiche bilan}{l'essentiel à retenir}\end{minipage}\par}

% Sommaire
\newtcolorbox{sommaire}{enhanced,colback=white,colframe=popink,boxrule=2.2pt,arc=10pt,
  drop shadow=popink,left=18pt,right=18pt,top=13pt,bottom=13pt,before skip=6pt,after skip=20pt,
  before upper={\poppill{Sommaire}{poppurple}{white}\par\vspace{9pt}\popsemi\fontsize{10.6}{14.5}\selectfont\color{popink}}}

% Signature de fin de cours — poussée tout en bas de la dernière page
\newcommand{\findecours}{%
  \par\vfill\nopagebreak
  \begin{center}\popsquiggle{poporange}\end{center}\vspace{4pt}
  \signatureonbuch
}
\AtBeginDocument{\def\llbracket{\mathopen{[\mkern-3.5mu[}}\def\rrbracket{\mathclose{]\mkern-3.5mu]}}} % intervalles d'entiers (glyphe ⟦ absent de la police)
% Glyphes absents de Fira Math : redéfinis à partir de glyphes disponibles
\AtBeginDocument{%
  \def\setminus{\mathbin{\backslash}}\let\smallsetminus\setminus
  \def\vdots{\mathord{\vbox{\baselineskip3.2pt\lineskiplimit0pt\kern4pt\hbox{.}\hbox{.}\hbox{.}}}}%
  \def\ddots{\mathinner{\mkern1mu\raise6.5pt\hbox{.}\mkern2mu\raise3.6pt\hbox{.}\mkern2mu\raise0.7pt\hbox{.}\mkern1mu}}%
  \def\triangle{\mathord{\scalebox{0.9}{$\symup{\Delta}$}}}%
  \def\nabla{\mathord{\raisebox{\depth}{\scalebox{1}[-1]{$\symup{\Delta}$}}}}%
  \def\checkmark{\popcheck{popdark}}%
  \def\star{\mathbin{\ast}}\def\bigstar{\mathord{\ast}}%
  \def\diamond{\mathbin{\scalebox{0.6}{$\Diamond$}}}%
  \def\bigcup{\mathop{\mathchoice{\raisebox{-0.3em}{\scalebox{1.7}{$\cup$}}}{\scalebox{1.25}{$\cup$}}{\cup}{\cup}}\displaylimits}%
  \def\bigcap{\mathop{\mathchoice{\raisebox{-0.3em}{\scalebox{1.7}{$\cap$}}}{\scalebox{1.25}{$\cap$}}{\cap}{\cap}}\displaylimits}%
  \def\lesseqqgtr{\mathrel{\lessgtr}}\def\bigoplus{\mathop{\oplus}\displaylimits}%
}
\providecommand{\ssi}{si et seulement si} % abréviation fréquente
\AtBeginDocument{\providecommand{\wideparen}{\overparen}} % arc de cercle

%% FIGFIX-BEGIN v3 — correctifs globaux des figures (docs/onbuch-generation/tools/figfix)
\usepackage{adjustbox}\usepackage{etoolbox}
% toute figure TikZ/pgfplots plus large que la ligne est réduite à la largeur disponible
\BeforeBeginEnvironment{tikzpicture}{\begin{adjustbox}{max width=\linewidth}}
\AfterEndEnvironment{tikzpicture}{\end{adjustbox}}
% légendes pgfplots lisibles par-dessus les courbes
\pgfplotsset{every axis/.append style={legend style={fill=white,fill opacity=0.92,text opacity=1,draw=black!15,inner sep=3pt}}}
% étiquettes d'arêtes (syntaxe edge["..."]) avec fond blanc
\tikzset{every edge quotes/.append style={fill=white,inner sep=1.2pt}}
% bulles de cartes mentales : texte à la ligne dans la bulle, bulle un peu plus grande
\tikzset{every concept/.append style={text width=2.0cm,align=center,minimum size=2.3cm,inner sep=2pt}}
%% FIGFIX-END
\begin{document}

\pagegarde{Lesson 43 — Reading: Texts about taking part in polls/surveys}{Anglais}{Tle A}{Chapter 4 : Citizenship and Human Rights}{EN43}{2 h}{In this reading lesson, students discover authentic-style texts about opinion polls and surveys, learn reading strategies (skimming, scanning, guessing meaning from context) and build the thematic vocabulary needed to understand and discuss surveys as a tool of citizenship.
\vspace{4pt}
\begin{multicols}{2}\small
\begin{itemize}
\item À la fin de cette leçon, je sais définir ce qu'est un sondage (poll) et une enquête (survey) et distinguer les deux notions.
\item Je sais utiliser le skimming pour dégager l'idée générale d'un texte en moins de deux minutes.
\item Je sais utiliser le scanning pour repérer rapidement une information précise (chiffre, date, nom).
\item Je sais déduire le sens d'un mot inconnu à partir du contexte et de la famille de mots.
\item Je sais répondre à des questions de compréhension variées (vrai/faux, QCM, réponses courtes) en justifiant par le texte.
\item Je connais le lexique essentiel du thème : respondent, sample, questionnaire, turnout, margin of error...
\end{itemize}\end{multicols}}

\begin{sommaire}
\begin{multicols}{2}
\begin{enumerate}
\item Warm-up: What is a poll?
\item Texte support : 'Why Polls Matter'
\item Compréhension : questions variées
\item Lexique thématique : the language of polls
\item Méthode de lecture : skimming, scanning, guessing
\item Production guidée : summary and opinion
\item Activité d'intégration
\item Méthodes et exercices résolus
\item Exercices
\item Corrigés des exercices
\item Fiche bilan
\end{enumerate}
\end{multicols}
\end{sommaire}

\begin{objectifs}
\begin{itemize}
\item À la fin de cette leçon, je sais définir ce qu'est un sondage (poll) et une enquête (survey) et distinguer les deux notions.
\item Je sais utiliser le skimming pour dégager l'idée générale d'un texte en moins de deux minutes.
\item Je sais utiliser le scanning pour repérer rapidement une information précise (chiffre, date, nom).
\item Je sais déduire le sens d'un mot inconnu à partir du contexte et de la famille de mots.
\item Je sais répondre à des questions de compréhension variées (vrai/faux, QCM, réponses courtes) en justifiant par le texte.
\item Je connais le lexique essentiel du thème : respondent, sample, questionnaire, turnout, margin of error...
\item Je sais résumer un texte sur les sondages en trois ou quatre phrases avec mes propres mots.
\item Je sais exprimer une opinion argumentée sur l'utilité des sondages dans la vie citoyenne.
\end{itemize}
\end{objectifs}

\begin{prerequis}
\begin{itemize}
\item Vocabulaire de base de la citoyenneté vu en début de chapitre (citizen, rights, vote, election).
\item Techniques de lecture rapide introduites en Première (repérage du titre, des mots-clés).
\item Questions en WH- et réponses courtes en anglais.
\item Expressions simples d'opinion : I think, I believe, in my opinion.
\item Nombres, pourcentages et dates en anglais.
\end{itemize}
\end{prerequis}

\newpage

\coursec{Warm-up: What is a poll?}

\courssub{Lead-in: a knock at the door in Ngoa-Ekelle}

\begin{textebox}[A real-life situation]
Last week, a team of researchers knocked at the door of a house in the Ngoa-Ekelle neighbourhood in Yaoundé: they were conducting a survey on mobile phone use among teenagers. Some residents answered willingly; others refused, suspicious. In the UK and the USA, opinion polls make the headlines every week — think of the Brexit referendum or American presidential polls. But what exactly is a poll? Who conducts surveys, and why should citizens take part? Today, we read texts about polls and surveys to understand this key tool of modern citizenship.
\end{textebox}

\textbf{Problématique.} \eng{What is a poll? Who conducts surveys, and why should citizens take part?} — Qu'est-ce qu'un sondage ? Qui mène les enquêtes, et pourquoi les citoyens devraient-ils y participer ? Gardez ces trois questions en tête pendant toute la leçon : le texte de la section suivante y répondra.

\courssub{Brainstorming: when someone asks you questions}

\begin{experience}[Class discussion — “Have you ever been surveyed?”]
Travail en binômes puis mise en commun au tableau. Répondez en anglais, en phrases complètes :
\begin{enumerate}
\item \eng{Has anyone ever asked you questions in the street, at the market, or on the phone? What about?}
\item \eng{Did you answer willingly or did you refuse? Why?}
\item \eng{Who do you think asks these questions: journalists, students, companies, the government?}
\item \eng{What do people do with your answers?}
\end{enumerate}
Exemple de production attendue : \eng{Last month, a student from the University of Yaoundé asked me questions about my favourite radio station. I answered because she was polite and it took only two minutes.} « Le mois dernier, une étudiante de l'université de Yaoundé m'a posé des questions sur ma station de radio préférée. J'ai répondu parce qu'elle était polie et que cela n'a pris que deux minutes. »
\end{experience}

Lors du brainstorming, les élèves réactivent naturellement du vocabulaire déjà connu. Fixons-le immédiatement au tableau sous forme de réseau :

\begin{center}
\begin{tabularx}{\linewidth}{|X|X|X|}
\hline
\rowcolor{popblueL} English & Français & Exemple en contexte \\
\hline
\cle{vote} (n./v.) & vote ; voter & \eng{People vote for their favourite candidate.} \\
\hline
\cle{election} (n.) & élection & \eng{The presidential election is next year.} \\
\hline
\cle{question} (n.) & question & \eng{The interviewer asked ten questions.} \\
\hline
\cle{answer} (n./v.) & réponse ; répondre & \eng{She gave a short answer.} \\
\hline
\cle{people} (n. pl.) & les gens, le peuple & \eng{Many people refused to answer.} \\
\hline
\cle{opinion} (n.) & opinion, avis & \eng{What is your opinion on school uniforms?} \\
\hline
\cle{to ask} (v.) & demander, poser (une question) & \eng{They asked me about my habits.} \\
\hline
\end{tabularx}
\end{center}

\begin{attention}[Faux amis et pièges de base]
\begin{itemize}
\item \eng{to ask a question} = poser une question. Ne dites jamais \emph{to pose a question} au niveau lycée (trop soutenu) et surtout pas \emph{to demand a question} : \eng{to demand} signifie « exiger ».
\item \eng{people} est déjà pluriel : \eng{People are...} et non \emph{peoples are}. \eng{Peoples} n'existe qu'au sens de « peuples » (ethnies, nations).
\item \eng{an answer \textbf{to} a question} : la préposition est \cle{to}, pas \emph{of}. \eng{the answer to question 3}.
\item \eng{to ask somebody a question} se construit \textbf{sans préposition} : \eng{She asked me a question}, et non \emph{she asked to me}.
\end{itemize}
\end{attention}

\courssub{Poll, survey or census? Three key words}

Avant de lire un texte sur les sondages, il faut distinguer trois mots que les journalistes anglophones emploient avec précision — et que le français rend souvent par le seul mot « sondage ».

\begin{definition}[Poll and survey]
A \cle{poll} is a study in which people are asked for their opinions about a subject or a person. A \cle{survey} is a more detailed investigation, usually with a questionnaire. A \cle{census} is an official count of all the people in a country.
\end{definition}

\begin{propriete}[Comment distinguer les trois notions]
\begin{itemize}
\item \textbf{Poll} (sondage d'opinion) : court, une à cinq questions, sur un échantillon (\eng{sample}) de personnes. Exemple typique : \eng{“Who will you vote for?”} On parle d'\eng{opinion poll}.
\item \textbf{Survey} (enquête) : plus large et plus détaillé, avec un \eng{questionnaire} complet (habitudes de vie, consommation, santé, éducation). Exemple : \eng{a survey on mobile phone use among teenagers}.
\item \textbf{Census} (recensement) : officiel, mené par l'État, il concerne \textbf{toute} la population, pas un échantillon. Au Cameroun comme au Royaume-Uni, le recensement a lieu environ tous les dix ans.
\end{itemize}
\end{propriete}

\begin{center}
\begin{tabularx}{\linewidth}{|X|X|X|X|}
\hline
\rowcolor{popblueL} Terme & Français & Taille / durée & Exemple anglais \\
\hline
\cle{poll} & sondage (d'opinion) & petit échantillon, quelques questions & \eng{An opinion poll predicts the next president.} \\
\hline
\cle{survey} & enquête, étude & questionnaire détaillé & \eng{A national survey on reading habits.} \\
\hline
\cle{census} & recensement & toute la population & \eng{The census counts every household.} \\
\hline
\cle{questionnaire} & questionnaire & liste de questions écrites & \eng{Fill in the questionnaire, please.} \\
\hline
\cle{sample} & échantillon & partie représentative & \eng{A sample of 1,000 voters was interviewed.} \\
\hline
\end{tabularx}
\end{center}

\begin{attention}[Homophones : poll / pole — et le double sens de “the polls”]
Ne pas confondre \eng{poll} (sondage) et \eng{pole} (poteau) — même prononciation « pôle », sens totalement différents : \eng{a telephone pole} (un poteau téléphonique), \eng{the North Pole} (le pôle Nord). De plus, \eng{the polls} au pluriel peut désigner les \textbf{bureaux de vote} : \eng{Voters went to the polls.} « Les électeurs se sont rendus aux urnes. » Le contexte décide du sens : \eng{The polls show a close race} (les sondages) vs \eng{The polls close at 6 p.m.} (les bureaux de vote).
\end{attention}

\begin{savaistu}[A short history of polling]
The first modern opinion poll was conducted in the USA in 1824 to predict a presidential election. Since then, polling has become a science: institutes choose a representative \eng{sample}, write neutral questions and publish a \eng{margin of error} (marge d'erreur). Today, Ghana and Nigeria have famous polling institutes, and Cameroonian media regularly publish surveys before elections. Polling is now part of democratic life on the whole continent.
\end{savaistu}

\courssub{Mind map: building the lexical field}

Voici la carte mentale du champ lexical que nous allons compléter pendant toute la leçon. Recopiez-la dans votre cahier et ajoutez un mot à chaque branche au fur et à mesure.

\begin{popfigure}
\begin{center}
\begin{tikzpicture}[
  mindmap,
  grow cyclic,
  every node/.style={concept, align=center, font=\small},
  concept color=popblue,
  level 1 concept/.append style={sibling angle=90, level distance=3.6cm},
  level 2 concept/.append style={sibling angle=60, level distance=2.6cm},
  scale=0.92, transform shape
]
\node[concept color=popblue, text=white, font=\bfseries\small, align=center]{POLLS \&\\SURVEYS}
  child[concept color=poporange] { node[align=center]{WHO?\\pollsters}
    child { node[concept color=poporangeL, text=popdark] {interviewer} }
    child { node[concept color=poporangeL, text=popdark] {respondent} }
  }
  child[concept color=popgreen] { node[align=center]{HOW?\\questionnaire}
    child { node[concept color=popgreenL, text=popdark] {by phone} }
    child { node[concept color=popgreenL, text=popdark] {online} }
  }
  child[concept color=poppurple] { node[align=center]{WHY?\\opinions}
    child { node[concept color=poppurpleL, text=popdark] {decisions} }
    child { node[concept color=poppurpleL, text=popdark] {predictions} }
  }
  child[concept color=poppink] { node[align=center]{WHERE?\\street}
    child { node[concept color=poppinkL, text=popdark] {school} }
    child { node[concept color=poppinkL, text=popdark] {media} }
  };
\end{tikzpicture}
\end{center}
\legende{Carte mentale du champ lexical : qui mène les sondages, comment, pourquoi et où.}
\end{popfigure}

\begin{definition}[Key words of the mind map]
\begin{itemize}
\item \cle{pollster} (n.) : sondeur, institut de sondage — la personne ou l'organisme qui mène le sondage.
\item \cle{respondent} (n.) : personne interrogée, répondant(e). On dit aussi \eng{interviewee}.
\item \cle{questionnaire} (n.) : liste écrite de questions (mot d'origine française, prononcé « kouèss-tchion-èr »).
\item \cle{sample} (n.) : échantillon, partie représentative de la population.
\end{itemize}
\end{definition}

\courssub{Predicting from the title and the picture}

\begin{methode}[Prédire avant de lire : une stratégie de bons lecteurs]
Avant de lire le moindre texte en anglais, suivez toujours ces quatre étapes :
\begin{enumerate}
\item \textbf{Lisez le titre.} Le texte de la prochaine section s'intitule \eng{“Why Polls Matter”}. Le verbe \eng{to matter} signifie « compter, avoir de l'importance ». Le titre annonce donc un texte qui \emph{justifie} l'importance des sondages.
\item \textbf{Regardez l'illustration.} Une image de personnes répondant à un questionnaire dans la rue confirme le thème : la collecte d'opinions.
\item \textbf{Formulez des hypothèses.} Posez-vous la question : \eng{What do I expect the text to be about?} Réponses possibles : \eng{I expect the text to explain why polls are useful for democracy.} / \eng{Maybe it will describe how pollsters work.}
\item \textbf{Anticipez le vocabulaire.} Listez mentalement les mots que le texte contiendra probablement : \eng{vote, opinion, sample, question, election, result, percentage}.
\end{enumerate}
\end{methode}

\begin{exemplebox}[Predictions made by Terminale A students]
\eng{“I think the text will explain why polls are important before elections.”} « Je pense que le texte expliquera pourquoi les sondages sont importants avant les élections. »

\eng{“In my opinion, the text will also warn us: polls can be wrong.”} « À mon avis, le texte nous mettra aussi en garde : les sondages peuvent se tromper. »

\eng{“I expect to read about the Brexit referendum, because the pollsters were surprised by the result.”} « Je m'attends à lire des choses sur le référendum du Brexit, car les sondeurs ont été surpris par le résultat. »
\end{exemplebox}

\courssub{A short warm-up text: “One question, one minute”}

Pour vérifier vos prédictions et activer le lexique, lisons d'abord ce court texte d'échauffement, plus simple que le texte principal de la leçon.

\begin{textebox}[One question, one minute]
Every Saturday morning, Mrs Etoundi sells tomatoes at the Mokolo market in Yaoundé. Last Saturday, a young man with a clipboard stopped in front of her stall. “Good morning, Madam. I work for a research institute. We are conducting a poll on food prices. May I ask you three questions? It will take one minute.”

Mrs Etoundi smiled. She is used to customers, but not to pollsters! “Why not?” she answered. The young man asked: “Have food prices increased this year? Do you think the government should control prices? And finally, do you trust opinion polls?”

Mrs Etoundi answered the first two questions easily, but the last one made her think. “I don't know,” she said. “Who reads your results?” “Newspapers, radio stations, and sometimes ministers,” the pollster replied. “Your opinion, together with the opinions of one thousand other people, helps leaders understand what citizens think.”

When he left, Mrs Etoundi felt proud. For the first time, somebody had asked for her opinion — and somebody was going to publish it. That evening, she told her daughter: “Today, I took part in a poll. My voice counts.”

Her daughter, a student in Terminale, smiled: “Mum, you have just discovered democracy's favourite tool!”
\end{textebox}

\textbf{Glossaire.}

\begin{center}
\begin{tabularx}{\linewidth}{|X|X|X|}
\hline
\rowcolor{popblueL} Mot anglais & Nature & Traduction \\
\hline
clipboard & n. & porte-bloc \\
\hline
stall & n. & étal, stand (au marché) \\
\hline
to conduct a poll & expr. v. & mener un sondage \\
\hline
to increase & v. & augmenter \\
\hline
to trust & v. & faire confiance à \\
\hline
pollster & n. & sondeur \\
\hline
citizen & n. & citoyen(ne) \\
\hline
my voice counts & expr. & ma voix compte \\
\hline
\end{tabularx}
\end{center}

\textbf{Questions de pré-compréhension (répondez oralement).}
\begin{enumerate}
\item \eng{Who stopped in front of Mrs Etoundi's stall?} (repérage d'information explicite)
\item \eng{How many questions did the pollster want to ask?} (scanning : cherchez un chiffre)
\item \eng{Why did Mrs Etoundi feel proud at the end?} (interprétation)
\item \eng{What does the daughter mean by “democracy's favourite tool”?} (déduction du sens)
\end{enumerate}

\courssub{Reading objectives: read fast, then read carefully}

\begin{aretenir}[Les deux lectures d'un texte]
Un bon lecteur ne lit jamais un texte une seule fois de la même manière :
\begin{itemize}
\item \textbf{Première lecture — \cle{skimming}} (lecture rapide, « survol ») : lire vite pour saisir l'\textbf{idée générale} (\eng{the gist}). On lit le titre, la première phrase de chaque paragraphe, la conclusion. Objectif : répondre à \eng{What is the text about?}
\item \textbf{Deuxième lecture — \cle{scanning}} (lecture ciblée) : relire en cherchant des \textbf{informations précises} : un chiffre, un nom, une date, une opinion. L'œil « balaie » le texte comme un scanner.
\item \textbf{En permanence — \cle{guessing}} : déduire le sens des mots inconnus grâce au contexte, sans dictionnaire. Exemple : dans \eng{a young man with a clipboard}, le contexte (marché, questions, institut) permet de deviner qu'un \eng{clipboard} sert à écrire.
\end{itemize}
\end{aretenir}

\begin{exoresolu}[Applying the strategy to the warm-up text]
Énoncé — Sans relire le texte en entier, répondez en une phrase : \eng{What is the text about?} (skimming). Puis trouvez en moins de trente secondes : \eng{How many people take part in the poll, according to the pollster?} (scanning).
\tcblower
Solution détaillée —
\textbf{1. Skimming.} Le titre (\eng{One question, one minute}), la première phrase (une vendeuse au marché Mokolo) et la dernière phrase (\eng{democracy's favourite tool}) suffisent : \eng{The text is about a market woman in Yaoundé who takes part in an opinion poll for the first time and discovers that her opinion matters.} « Le texte parle d'une vendeuse de Yaoundé qui participe pour la première fois à un sondage et découvre que son opinion compte. »
\textbf{2. Scanning.} On cherche un chiffre lié aux participants. Au troisième paragraphe : \eng{“the opinions of one thousand other people”}. Réponse : \eng{About one thousand people take part in the poll.} Piège : ne pas confondre avec les \emph{trois questions} posées à Mrs Etoundi — le scanning exige de lire la question avec précision.
\end{exoresolu}

\begin{exercice}[Your turn — before the main text]
\begin{enumerate}
\item Traduisez en anglais : « Les sondeurs posent des questions aux citoyens dans la rue. »
\item Complétez avec \eng{poll}, \eng{survey} ou \eng{census} : (a) \eng{The government organises a \ldots\ every ten years to count the population.} (b) \eng{A quick \ldots\ shows that most students prefer English in the morning.} (c) \eng{The university published a long \ldots\ on students' eating habits, with forty questions.}
\item Devinez le sens de \eng{clipboard} et de \eng{stall} sans dictionnaire, puis vérifiez avec le glossaire.
\item Formulez votre propre prédiction sur le texte \eng{“Why Polls Matter”} en une phrase commençant par \eng{I expect the text to\ldots}
\end{enumerate}
\end{exercice}

\begin{corrige}[Exercice « Your turn »]
\begin{enumerate}
\item \eng{Pollsters ask citizens questions in the street.} (Attention : \eng{to ask somebody questions}, sans préposition ; \eng{citizens} = citoyens.)
\item (a) \eng{census} — toute la population, tous les dix ans ; (b) \eng{poll} — rapide, une seule question d'opinion ; (c) \eng{survey} — enquête longue et détaillée (quarante questions).
\item \eng{clipboard} = porte-bloc (contexte : le jeune homme écrit les réponses) ; \eng{stall} = étal (contexte : Mrs Etoundi y vend des tomates).
\item Exemple : \eng{I expect the text to explain why polls are important for citizens and for democracy.}
\end{enumerate}
\end{corrige}

\begin{aretenir}[To sum up — what you must remember]
\begin{itemize}
\item A \cle{poll} is short and measures opinions ; a \cle{survey} is a detailed investigation with a questionnaire ; a \cle{census} counts the whole population.
\item \eng{Poll} et \eng{pole} se prononcent pareil (« pôle ») mais ne veulent pas dire la même chose ; \eng{the polls} peut aussi signifier « les bureaux de vote ».
\item Avant de lire : prédire à partir du titre et de l'illustration (\eng{What do I expect?}).
\item Pendant la lecture : \cle{skimming} pour l'idée générale, \cle{scanning} pour les détails, \cle{guessing} pour les mots inconnus.
\end{itemize}
Vous êtes maintenant prêts pour le texte principal : \eng{“Why Polls Matter”}.
\end{aretenir}

\coursec{Texte support : ‘Why Polls Matter’}

Dans la section précédente, nous avons découvert ce qu'est un \cle{poll} (un sondage) et pourquoi les citoyens y participent. Nous allons maintenant lire un article journalistique complet sur ce thème. C'est le cœur de la leçon : vous allez apprendre à lire un texte authentique en deux étapes — d'abord une lecture globale (\cle{skimming}), puis une lecture détaillée (\cle{scanning}) — et à deviner le sens des mots inconnus grâce au contexte avant de vérifier dans le glossaire. Ces trois stratégies sont exactement celles qu'exige l'épreuve de compréhension de l'écrit au Baccalauréat.

\courssub{Lecture silencieuse : le texte}

\begin{methode}[Avant de lire : trois réflexes]
\begin{enumerate}
\item Regardez le \textbf{titre} : \eng{Why Polls Matter}. Que vous annonce-t-il ? Le verbe \eng{to matter} signifie « compter, avoir de l'importance ». Le titre promet donc des \emph{arguments} expliquant l'importance des sondages.
\item Regardez la \textbf{longueur} et la \textbf{mise en page} : un seul texte continu en cinq paragraphes, style article de presse, environ 250 mots.
\item Fixez-vous un \textbf{objectif de lecture} : « De quoi parle ce texte ? L'auteur est-il pour ou contre les sondages ? »
\end{enumerate}
Lisez maintenant le texte \textbf{silencieusement et sans dictionnaire}, en deux minutes maximum.
\end{methode}

\begin{textebox}[Why Polls Matter — a press article]
Every election year, pollsters telephone hundreds of people in Yaoundé, Douala and Bamenda. They ask simple questions: “Who will you vote for? Do you trust the government? Are you satisfied with the roads in your region?” These surveys use a sample — a small group of people chosen to represent the whole population. If the sample is well chosen, the results reflect the opinion of millions of citizens who were never telephoned at all.

Conducting a serious poll is not easy. Pollsters must write clear questions, contact respondents in towns and in villages, and record the answers honestly. Then they publish their findings: for example, “sixty per cent of respondents say that education is their main concern.”

In Britain, polls predicted the result of the 2016 Brexit referendum with a margin of error of only three per cent. This means the real figures could be three points higher or lower than the published ones. In Cameroon, polls also try to estimate voter turnout before election day, so that journalists and observers know what to expect.

Critics say that polls can influence voters instead of simply measuring opinion: if people read that one candidate is far ahead, some of them may not bother to vote. Supporters answer that polls give a voice to ordinary citizens between elections, because leaders can see what the population really thinks.

Whether we like them or not, polls have become an essential tool of modern democracy — and taking part in a poll is a small but real act of citizenship.
\end{textebox}

\textbf{Glossaire intégré} (à consulter \emph{après} avoir fait vos hypothèses de sens) :

\begin{center}
\begin{tabularx}{\linewidth}{|l|l|X|}
\hline
\rowcolor{popblueL} English word & Nature & Français \\
\hline
pollster & noun & sondeur, institut de sondage \\
\hline
survey & noun & enquête, sondage \\
\hline
sample & noun & échantillon \\
\hline
respondent & noun & personne interrogée, répondant \\
\hline
findings & noun (plur.) & résultats, conclusions d'une enquête \\
\hline
margin of error & noun phrase & marge d'erreur \\
\hline
turnout & noun & taux de participation \\
\hline
to conduct (a poll) & verb & mener, réaliser (un sondage) \\
\hline
to predict & verb & prédire \\
\hline
to bother to vote & verb phrase & se donner la peine de voter \\
\hline
\end{tabularx}
\end{center}

\courssub{Première lecture globale : skimming}

Le \cle{skimming} (lecture rapide, « survol ») consiste à lire vite pour saisir le \textbf{sens général}, sans s'arrêter sur les mots inconnus. On lit surtout le titre, la première phrase de chaque paragraphe et la conclusion.

\begin{exemplebox}[Ce que le skimming doit vous donner en 2 minutes]
\begin{itemize}
\item \textbf{Thème (topic)} : the role of opinion polls in a democracy. (Le rôle des sondages en démocratie.)
\item \textbf{Type de texte} : a press article / an opinion piece — un article de presse argumentatif, avec des exemples du Cameroun et de la Grande-Bretagne.
\item \textbf{Point de vue de l'auteur} : balanced but positive — il présente les critiques (\eng{critics say...}) et les défenseurs (\eng{supporters answer...}), puis conclut que les sondages sont \eng{an essential tool of modern democracy}.
\end{itemize}
\end{exemplebox}

\begin{aretenir}[La question-clé du skimming]
Après un survol, vous devez pouvoir répondre en une phrase : \eng{What is the text about?} — \eng{The text is about why polls are important in a democracy.} Si vous ne pouvez pas formuler cette phrase, relisez le titre et la première phrase de chaque paragraphe.
\end{aretenir}

\courssub{Deuxième lecture détaillée : scanning}

Le \cle{scanning} (lecture sélective) consiste à chercher une \textbf{information précise} : un chiffre, un lieu, un exemple, un argument. L'œil court sur le texte sans tout lire, comme quand on cherche un nom dans un annuaire.

\begin{exemplebox}[Repérage dans ‘Why Polls Matter’]
\begin{itemize}
\item \textbf{Chiffres} : \eng{sixty per cent} (main concern : education) ; \eng{three per cent} (margin of error, Brexit referendum) ; \eng{2016} (date du référendum).
\item \textbf{Lieux} : Yaoundé, Douala, Bamenda (Cameroun) ; Britain (Grande-Bretagne).
\item \textbf{Exemple camerounais} : les sondages téléphoniques et l'estimation du \eng{voter turnout} avant le jour du scrutin.
\item \textbf{Exemple britannique} : les sondages ont prédit le résultat du référendum sur le Brexit avec une faible marge d'erreur.
\item \textbf{Argument contre} : polls can \emph{influence} voters instead of measuring opinion.
\item \textbf{Argument pour} : polls give a \emph{voice} to ordinary citizens between elections.
\end{itemize}
\end{exemplebox}

\begin{popfigure}
\begin{center}
\begin{tikzpicture}[node distance=0.45cm]
\node[draw=popblue, fill=popblueL, rounded corners, minimum width=2.3cm, minimum height=0.9cm, align=center] (q) {\textbf{question}};
\node[draw=poppurple, fill=poppurpleL, rounded corners, minimum width=2.6cm, minimum height=0.9cm, align=center, right=of q] (s) {\textbf{sample}\\ of people};
\node[draw=poporange, fill=poporangeL, rounded corners, minimum width=2.3cm, minimum height=0.9cm, align=center, right=of s] (a) {\textbf{answers}};
\node[draw=popgreen, fill=popgreenL, rounded corners, minimum width=2.3cm, minimum height=0.9cm, align=center, right=of a] (r) {\textbf{results}};
\node[draw=poppink, fill=poppinkL, rounded corners, minimum width=2.5cm, minimum height=0.9cm, align=center, right=of r] (d) {\textbf{decisions}};
\draw[-{Stealth}, very thick, popdark] (q) -- (s);
\draw[-{Stealth}, very thick, popdark] (s) -- (a);
\draw[-{Stealth}, very thick, popdark] (a) -- (r);
\draw[-{Stealth}, very thick, popdark] (r) -- (d);
\end{tikzpicture}
\end{center}
\legende{How a poll works : question $\to$ sample of people $\to$ answers $\to$ results $\to$ decisions}
\end{popfigure}

\courssub{Deviner le sens des mots inconnus (guessing from context)}

Avant de consulter le glossaire, un bon lecteur formule des \textbf{hypothèses}. Trois indices vous aident :

\begin{enumerate}
\item \textbf{La nature du mot} : dans \eng{pollsters telephone hundreds of people}, \eng{pollsters} est un nom pluriel, sujet du verbe — ce sont donc des \emph{personnes} (celles qui téléphonent).
\item \textbf{Le contexte logique} : dans \eng{a sample — a small group chosen to represent the whole population}, le tiret introduit une \emph{définition} : un échantillon.
\item \textbf{La famille de mots} : \eng{respondent} contient \eng{respond} (répondre) + le suffixe \eng{-ent} qui désigne une personne (comme \eng{student}, \eng{assistant}) : la personne qui répond.
\end{enumerate}

\begin{definition}[Margin of error]
\eng{The margin of error is the amount by which poll results may differ from reality.} — La \cle{marge d'erreur} est l'écart possible entre les résultats publiés d'un sondage et la réalité. \eng{A margin of error of three per cent} signifie que le chiffre réel peut être trois points au-dessus ou au-dessous du chiffre annoncé.
\end{definition}

\begin{exemplebox}[Exemples traduits]
\eng{The survey was conducted among 1,000 respondents.} — « L'enquête a été menée auprès de 1 000 personnes interrogées. »

\eng{Voter turnout was low.} — « Le taux de participation électorale était faible. »

\eng{The pollsters published their findings in the newspaper.} — « Les sondeurs ont publié leurs résultats dans le journal. »

\eng{The sample must represent the whole population.} — « L'échantillon doit représenter l'ensemble de la population. »
\end{exemplebox}

\begin{attention}[Faux amis et constructions piégeuses]
\begin{itemize}
\item \eng{to take part in a poll} = participer à un sondage. Ne dites jamais \emph{to participate to a poll} : \eng{participate} se construit avec la préposition \cle{in}, comme \eng{to take part in}.
\item \eng{to conduct} signifie « mener, réaliser » (une enquête) ; le « conducteur » d'un bus se dit \eng{driver}.
\item \eng{findings} (résultats d'enquête) s'emploie presque toujours au \textbf{pluriel} ; on dira \eng{the findings show that...} et non \emph{a finding shows}.
\item \eng{actually} ne veut pas dire « actuellement » mais « en fait » ; « actuellement » se dit \eng{currently} ou \eng{at present}.
\item Attention à la prononciation de \eng{survey} : le nom se prononce « seu-veï » (accent sur la première syllabe), le verbe \eng{to survey} « se-veï » (accent sur la deuxième).
\end{itemize}
\end{attention}

\courssub{Tableau récapitulatif : le vocabulaire du sondage}

\begin{center}
\begin{tabularx}{\linewidth}{|X|X|X|}
\hline
\rowcolor{popblueL} English & Français & Exemple en contexte \\
\hline
poll / survey & sondage / enquête & \eng{They conducted a survey in Douala.} \\
\hline
pollster & sondeur & \eng{Pollsters telephoned 500 people.} \\
\hline
sample & échantillon & \eng{The sample represents the population.} \\
\hline
respondent & personne interrogée & \eng{Each respondent answered five questions.} \\
\hline
findings & résultats & \eng{The findings were published on Monday.} \\
\hline
margin of error & marge d'erreur & \eng{A margin of error of three per cent.} \\
\hline
turnout & taux de participation & \eng{Turnout was high among young voters.} \\
\hline
to take part in & participer à & \eng{She took part in the poll.} \\
\hline
\end{tabularx}
\end{center}

\courssub{Exercice résolu et entraînement}

\begin{exoresolu}[Guessing meaning from context]
Énoncé — Sans dictionnaire, devinez le sens du mot en gras dans cette phrase tirée du texte : \eng{“Supporters answer that polls give a \textbf{voice} to ordinary citizens between elections.”} Puis traduisez la phrase entière.
\tcblower
Solution détaillée — Étape 1 : nature du mot. \eng{Voice} est un nom, objet du verbe \eng{give} ; on « donne » quelque chose aux citoyens. Étape 2 : contexte. La phrase oppose les sondages aux élections : entre deux élections, les citoyens ne votent pas, donc ils ne peuvent pas s'exprimer... sauf grâce aux sondages. Étape 3 : hypothèse. \eng{Voice} = « voix », au sens figuré : une façon de s'exprimer, de se faire entendre. Traduction complète : « Les partisans répondent que les sondages donnent une voix aux citoyens ordinaires entre les élections. » Vérification : le sens figuré de \eng{voice} existe aussi en français (« avoir voix au chapitre »), ce qui confirme l'hypothèse.
\end{exoresolu}

\begin{exercice}[Skimming and scanning practice]
\begin{enumerate}
\item (Skimming) In one sentence, say what the text \eng{Why Polls Matter} is about.
\item (Scanning) Find in the text : (a) the margin of error of the British polls in 2016 ; (b) the three Cameroonian cities mentioned ; (c) the main concern of sixty per cent of respondents.
\item (Guessing) Guess the meaning of : \eng{to trust}, \eng{far ahead}, \eng{whether we like them or not}.
\item (Grammar watch) Rewrite with the correct preposition : \eng{Many citizens participate to polls every year.}
\end{enumerate}
\end{exercice}

\begin{corrige}[Exercice : Skimming and scanning practice]
\begin{enumerate}
\item \eng{The text is about the importance of opinion polls in modern democracy, with examples from Cameroon and Britain.}
\item (a) \eng{Three per cent.} (b) \eng{Yaoundé, Douala and Bamenda.} (c) \eng{Education.}
\item \eng{to trust} = avoir confiance en ; \eng{far ahead} = loin en tête ; \eng{whether we like them or not} = que nous les aimions ou non.
\item \eng{Many citizens participate \textbf{in} polls every year.} (ou mieux : \eng{take part in polls}).
\end{enumerate}
\end{corrige}

\begin{savaistu}[Did you know?]
Le mot anglais \eng{poll} vient d'un vieux mot germanique qui signifiait « tête » : à l'origine, faire un \eng{poll}, c'était littéralement « compter les têtes » lors d'un vote ! C'est aussi pourquoi le lieu où l'on vote s'appelle \eng{polling station} (bureau de vote) et le jour du scrutin \eng{polling day}. Au Cameroun anglophone, on entend souvent \eng{opinion poll} à la radio pendant les périodes électorales — un bel exemple de mot dont l'histoire éclaire le sens actuel.
\end{savaistu}

Vous savez maintenant lire ce texte en deux passages complémentaires : le \cle{skimming} pour l'idée générale, le \cle{scanning} pour les détails précis, et vous pouvez deviner le sens des mots inconnus grâce au contexte. La section suivante mettra ces stratégies à l'épreuve avec des questions de compréhension et d'interprétation variées, comme à l'examen.

\coursec{Compréhension : questions variées}

Dans la section précédente, nous avons découvert le texte support \emph{Why Polls Matter} et nous avons effectué une première lecture globale (\cle{skimming}). Il est maintenant temps de passer à la \cle{lecture approfondie} (\eng{detailed reading}) : c'est l'étape où l'on vérifie, ligne par ligne, que l'on a vraiment compris ce que l'auteur dit — et ce qu'il ne dit pas. Au baccalauréat, cette étape correspond aux \eng{reading comprehension questions}, qui représentent souvent la moitié des points de l'épreuve de compréhension de l'écrit. Nous allons donc nous entraîner sur les cinq grands types de questions : le vrai/faux, le QCM, les questions à réponses courtes, les questions d'interprétation et la recherche de mots dans le texte.

\courssub{Rappel du texte support}

Pour travailler efficacement, relisons le texte en y repérant les lignes : au Bac, on vous demandera souvent de justifier vos réponses par une citation précise (\eng{line reference}).

\begin{textebox}[Why Polls Matter (rappel du texte support)]
Every week, somewhere in the world, a polling company publishes new figures: who is leading the presidential race, what citizens think about the economy, or how many young people trust the media. But what exactly is a poll, and why should we care about the results? (lines 1--4)

A poll, or survey, is a study in which a small group of people, called a sample, is asked questions in order to understand the opinions of a much larger population. A sample is a small group of people chosen to represent the whole population. If the sample is well chosen, a thousand interviews can give a reliable picture of what millions of citizens think. Polls are conducted by specialised companies, by universities, and sometimes by newspapers or television channels. (lines 5--11)

Polls matter because they give a voice to ordinary citizens. Between two elections, they are often the only way for people to tell leaders what they really think about schools, hospitals or roads. In Britain, for example, opinion polls before the 2016 referendum revealed how divided the country was about leaving the European Union. (lines 12--17)

However, polls are not perfect. A badly chosen sample, or a badly written question, can produce misleading results. That is why readers should always ask: who paid for this poll, how many people were interviewed, and how was the question asked? A critical reader is a protected citizen. (lines 18--22)
\end{textebox}

\textbf{Glossaire :}
\begin{itemize}
\item \eng{poll} (n.) : sondage ; \eng{survey} (n.) : enquête, étude.
\item \eng{sample} (n.) : échantillon ; \eng{population} (n.) : population (ensemble des personnes concernées).
\item \eng{to conduct a poll} (v.) : réaliser un sondage ; \eng{reliable} (adj.) : fiable.
\item \eng{to give a voice to} (expr.) : donner la parole à ; \eng{misleading} (adj.) : trompeur.
\item \eng{referendum} (n.) : référendum ; \eng{divided} (adj.) : divisé.
\end{itemize}

\courssub{Méthode : répondre à une question de compréhension}

\begin{methode}[Les quatre étapes d'une bonne réponse]
\begin{enumerate}
\item \textbf{Lire la question avant le texte} (ou la relire avant de chercher) : elle vous dit exactement quoi chercher.
\item \textbf{Souligner les mots-clés de la question} (\eng{key words}) : noms propres, verbes, notions (\eng{Who conducts polls?} → mots-clés : \eng{who}, \eng{conduct}).
\item \textbf{Repérer le passage correspondant} dans le texte grâce à ces mots-clés ou à leurs synonymes (\eng{conduct} = \eng{carry out}).
\item \textbf{Répondre par une phrase complète}, avec un sujet et un verbe, en citant le texte ou en le reformulant, et en indiquant la ligne (\eng{line 8}).
\end{enumerate}
\end{methode}

\begin{exemplebox}[Une question, une réponse modèle]
\textbf{Question :} \eng{What is a sample?}

\textbf{Réponse attendue :} \eng{A sample is a small group of people chosen to represent the whole population.} (lines 6--7)

Traduction : « Un échantillon est un petit groupe de personnes choisi pour représenter l'ensemble de la population. »

Remarquez la structure de la réponse : elle reprend les mots de la question (\eng{a sample is...}) et ajoute l'information du texte. Une réponse comme \eng{small group} seul est incomplète : elle perd des points.
\end{exemplebox}

\courssub{Type 1 : les questions Vrai / Faux (True or False)}

Le vrai/faux teste votre compréhension des \textbf{détails} du texte. La règle d'or : une affirmation n'est vraie que si le texte la confirme \emph{explicitement}. Si le texte ne dit rien, ou dit le contraire, la réponse est \eng{False} (ou \eng{Not stated} si cette option existe).

\begin{exoresolu}[True or False — justifiez par une ligne du texte]
\textbf{Énoncé.} Say whether the following statements are True or False. Justify each answer with a line from the text.

a) \eng{Polls only exist in Britain.}

b) \eng{A well-chosen sample of a thousand people can reflect the opinion of millions.}

c) \eng{Polls are always accurate.}
\tcblower
\textbf{Solution détaillée.}

a) \textbf{False.} Le texte dit \eng{Every week, somewhere in the world, a polling company publishes new figures} (line 1) : les sondages existent partout dans le monde, pas seulement en Grande-Bretagne. La Grande-Bretagne n'est citée que comme \emph{exemple} (line 15). Piège classique : transformer un exemple en généralité.

b) \textbf{True.} \eng{If the sample is well chosen, a thousand interviews can give a reliable picture of what millions of citizens think} (lines 8--10). L'affirmation reprend presque mot pour mot cette phrase.

c) \textbf{False.} \eng{However, polls are not perfect... can produce misleading results} (lines 18--20). Le mot \eng{always} est un signal d'alarme : les affirmations absolues (\eng{always}, \eng{never}, \eng{only}, \eng{all}) sont souvent fausses.
\end{exoresolu}

\courssub{Type 2 : le QCM (idée principale et ton de l'auteur)}

Le QCM teste la \textbf{compréhension globale} : l'idée principale du texte (\eng{main idea}) et le ton de l'auteur (\eng{the author's tone}). Pour le ton, trois options reviennent souvent :

\begin{center}
\begin{tabularx}{\linewidth}{|X|X|X|}
\hline
\rowcolor{popblueL} Ton & English & Indices dans le texte \\
\hline
Neutre & neutral & faits, définitions, pas de jugement \\
\hline
Critique & critical & \eng{however}, \eng{not perfect}, mises en garde \\
\hline
Enthousiaste & enthusiastic & adjectifs élogieux, exclamations \\
\hline
\end{tabularx}
\end{center}

\begin{exoresolu}[QCM — main idea and tone]
\textbf{Énoncé.} Choose the best answer.

1) \eng{The main idea of the text is:}
a) \eng{Polls are useless and dangerous.}
b) \eng{Polls are useful tools for democracy, but they must be read critically.}
c) \eng{British polls are the best in the world.}

2) \eng{The author's tone is:}
a) \eng{enthusiastic} \quad b) \eng{aggressive} \quad c) \eng{balanced and critical}
\tcblower
\textbf{Solution détaillée.}

1) \textbf{b)} est la bonne réponse. Le texte présente d'abord l'utilité des sondages (\eng{they give a voice to ordinary citizens}, line 13), puis leurs limites (\eng{polls are not perfect}, line 18). Les options a) et c) sont des extrêmes : le texte ne dit ni que les sondages sont inutiles, ni que les sondages britanniques sont les meilleurs.

2) \textbf{c)} est la bonne réponse. L'auteur reconnaît la valeur des sondages mais invite à la prudence (\eng{A critical reader is a protected citizen}, line 22) : le ton est donc équilibré (\eng{balanced}) et critique (\eng{critical}), jamais enthousiaste ni agressif.
\end{exoresolu}

\courssub{Type 3 : les questions à réponses courtes (short answers)}

Ces questions commencent par \eng{Who}, \eng{What}, \eng{Where}, \eng{When}, \eng{Why}, \eng{How many}. Elles testent le \cle{scanning} : repérer rapidement une information précise. Répondez par une phrase complète.

\begin{exemplebox}[Questions et réponses modèles]
\begin{itemize}
\item \eng{Who conducts polls?} → \eng{Polls are conducted by specialised companies, universities, and sometimes newspapers or television channels.} (lines 10--11) « Qui réalise les sondages ? »
\item \eng{What is a sample?} → \eng{A sample is a small group of people chosen to represent the whole population.} (lines 6--7) « Qu'est-ce qu'un échantillon ? »
\item \eng{What example from Britain is given?} → \eng{The text gives the example of the opinion polls before the 2016 referendum, which revealed how divided the country was about leaving the European Union.} (lines 15--17) « Quel exemple britannique est donné ? »
\end{itemize}
\end{exemplebox}

\courssub{Type 4 : la question d'interprétation}

La question d'interprétation exige de \textbf{lire entre les lignes} (\eng{to read between the lines}) : la réponse n'est pas écrite telle quelle, il faut la déduire.

\begin{exoresolu}[Question d'interprétation]
\textbf{Énoncé.} \eng{Why does the author say that polls give a voice to citizens?} (Answer in your own words.)
\tcblower
\textbf{Solution détaillée.} Le passage clé est aux lignes 12--14 : \eng{Between two elections, they are often the only way for people to tell leaders what they really think.} On en déduit la réponse, reformulée : \eng{The author means that citizens vote only every few years, so polls are their only regular opportunity to express their opinions to their leaders between elections.} « Les citoyens ne votent que tous les quelques années ; les sondages sont donc leur seule occasion régulière d'exprimer leur opinion aux dirigeants entre deux élections. » Notez l'usage de \eng{the author means that...} pour signaler qu'il s'agit d'une interprétation.
\end{exoresolu}

\courssub{Type 5 : trouver le mot dans le texte (find the word that means...)}

On vous donne une définition ; vous devez retrouver le mot du texte qui y correspond. Stratégie : cherchez un synonyme ou une paraphrase de la définition dans le texte.

\begin{center}
\begin{tabularx}{\linewidth}{|X|X|X|}
\hline
\rowcolor{popblueL} Définition donnée & Mot du texte & Ligne \\
\hline
\eng{a small group representing a larger population} & \eng{sample} & 6 \\
\hline
\eng{trustworthy, that you can depend on} & \eng{reliable} & 9 \\
\hline
\eng{giving a false idea, deceptive} & \eng{misleading} & 20 \\
\hline
\eng{a vote in which all citizens decide one question} & \eng{referendum} & 16 \\
\hline
\eng{carried out, organised} & \eng{conducted} & 10 \\
\hline
\end{tabularx}
\end{center}

\courssub{Synthèse : quel type de question teste quoi ?}

\begin{popfigure}
\begin{tikzpicture}[font=\small]
\node[draw=popblue, fill=popblueL, rounded corners, minimum width=5.2cm, minimum height=1.1cm, align=center] (q1) at (0,3.3) {\textbf{True / False}};
\node[draw=poporange, fill=poporangeL, rounded corners, minimum width=6.4cm, minimum height=1.1cm, align=center] (t1) at (7.2,3.3) {compréhension des \textbf{détails}\\ repérage + justification par la ligne};
\node[draw=popblue, fill=popblueL, rounded corners, minimum width=5.2cm, minimum height=1.1cm, align=center] (q2) at (0,1.9) {\textbf{QCM}};
\node[draw=poporange, fill=poporangeL, rounded corners, minimum width=6.4cm, minimum height=1.1cm, align=center] (t2) at (7.2,1.9) {\textbf{idée globale} et ton de l'auteur\\ (skimming + lecture critique)};
\node[draw=popblue, fill=popblueL, rounded corners, minimum width=5.2cm, minimum height=1.1cm, align=center] (q3) at (0,0.5) {\textbf{Short answers}};
\node[draw=poporange, fill=poporangeL, rounded corners, minimum width=6.4cm, minimum height=1.1cm, align=center] (t3) at (7.2,0.5) {\textbf{scanning} : information précise\\ (who, what, where, how many)};
\node[draw=popblue, fill=popblueL, rounded corners, minimum width=5.2cm, minimum height=1.1cm, align=center] (q4) at (0,-0.9) {\textbf{Interpretation}};
\node[draw=poporange, fill=poporangeL, rounded corners, minimum width=6.4cm, minimum height=1.1cm, align=center] (t4) at (7.2,-0.9) {\textbf{lecture entre les lignes}\\ déduction + reformulation};
\draw[-{Stealth}, thick, popmuted] (q1) -- (t1);
\draw[-{Stealth}, thick, popmuted] (q2) -- (t2);
\draw[-{Stealth}, thick, popmuted] (q3) -- (t3);
\draw[-{Stealth}, thick, popmuted] (q4) -- (t4);
\end{tikzpicture}
\legende{Chaque type de question teste une compétence de lecture différente.}
\end{popfigure}

\courssub{Correction collective : exiger la justification}

En classe, la correction se fait au tableau, et aucune réponse n'est acceptée sans sa \textbf{preuve}. L'enseignant demande systématiquement : \eng{Where did you find this? Which line?} L'élève doit répondre en citant le texte : \eng{It says in line 8 that...} ou \eng{According to the text (line 15),...}

\begin{aretenir}[Les expressions pour justifier une réponse]
\begin{itemize}
\item \eng{According to the text, ...} : « D'après le texte, ... »
\item \eng{The author says (in line 5) that ...} : « L'auteur dit (à la ligne 5) que ... »
\item \eng{This is shown by the example of ...} : « Cela est montré par l'exemple de ... »
\item \eng{The text does not say that...} : « Le texte ne dit pas que... » (pour un \eng{False})
\end{itemize}
\end{aretenir}

\begin{attention}[Le piège du recopiage au Bac]
Au baccalauréat, une réponse juste mais \textbf{recopiée mot pour mot sans rapport avec la question} ne vaut rien : le correcteur voit immédiatement que vous n'avez pas compris. Deux règles vitales :
\begin{enumerate}
\item Quand la consigne dit \eng{in your own words} (« avec vos propres mots »), vous devez \textbf{reformuler} : remplacez les mots du texte par des synonymes (\eng{reliable} → \eng{trustworthy} ; \eng{citizens} → \eng{ordinary people}).
\item Ne recopiez jamais un paragraphe entier en espérant que la réponse s'y trouve : citez uniquement le passage pertinent, ou reformulez-le.
\end{enumerate}
Autre piège fréquent des francophones : répondre en français ! Toute la réponse doit être en anglais, avec un sujet et un verbe conjugué. Évitez aussi les réponses sans verbe comme \eng{Yes, polls matter.} quand la question commence par \eng{Why} : il faut \eng{Polls matter because...}
\end{attention}

\begin{exercice}[À vous de jouer]
Répondez en anglais, par des phrases complètes, en indiquant la ligne qui justifie chaque réponse.

1) True or False: \eng{Only governments conduct polls.}

2) \eng{How many interviews can give a reliable picture of public opinion?}

3) \eng{Find the word in the text that means ``not telling the whole truth, likely to deceive''.}

4) \eng{In your own words, what three questions should a critical reader ask about a poll?}
\end{exercice}

\begin{corrige}[Exercice « À vous de jouer »]
1) \textbf{False.} \eng{Polls are conducted by specialised companies, by universities, and sometimes by newspapers or television channels} (lines 10--11) : les gouvernements ne sont même pas mentionnés.

2) \eng{A thousand interviews can give a reliable picture of what millions of citizens think, if the sample is well chosen.} (lines 8--10)

3) Le mot est \eng{misleading} (line 20) : « trompeur ».

4) Réponse modèle reformulée : \eng{A critical reader should ask who financed the poll, how many people answered it, and how the question was formulated.} (d'après les lignes 20--22, avec \eng{paid for} → \eng{financed}, \eng{interviewed} → \eng{answered}, \eng{asked} → \eng{formulated}).
\end{corrige}

\coursec{Lexique thématique : the language of polls}

Dans la section précédente, nous avons répondu à des questions de compréhension sur le texte \eng{Why Polls Matter}. Vous avez certainement remarqué que certaines difficultés venaient du \cle{vocabulaire} : mots comme \eng{respondents}, \eng{turnout}, \eng{findings} ou \eng{margin of error} appartiennent à un champ lexical très précis, celui des \cle{sondages et des enquêtes}. Cette section a pour but de vous faire maîtriser ce vocabulaire : vous devez pouvoir le reconnaître à la lecture, le prononcer correctement et l'employer dans vos propres phrases — une compétence directement évaluée au Baccalauréat, tant en compréhension de l'écrit qu'en composition. Nous allons classer ces mots en quatre familles logiques : les \cle{acteurs}, les \cle{outils}, les \cle{actions} et les \cle{résultats}.

\courssub{Observation : le vocabulaire dans son contexte}

Avant toute liste de mots, lisons ce court article de presse original. Soulignez mentalement tous les mots liés aux sondages : ce sont eux que nous allons étudier.

\begin{textebox}[Newspaper article — “Students Speak Out in National Survey”]
A team of researchers from the University of Buea has just published the findings of a nationwide survey on young people's reading habits. The pollsters interviewed a sample of two thousand students in secondary schools across Cameroon, from Yaoundé to Bamenda. Each respondent answered a short questionnaire of twenty questions, either on paper forms or online.

The results surprised many teachers. Sixty per cent of respondents said they read for pleasure at least once a week, while only twenty-five per cent said they never opened a book outside school. The majority of students — about seventy per cent — prefer reading on their phones. However, the researchers warn that the margin of error is three per cent, so the figures must be read carefully.

“We wanted reliable results, so we conducted the survey over three months,” explains Dr Etame, who led the research team. “The turnout was excellent: almost every student who was asked agreed to take part.” The team noticed a clear trend: the older the students are, the more they read news articles rather than novels. The full findings will be published next month in an education journal, and the Ministry of Secondary Education has already asked to see them.
\end{textebox}

\textbf{Glossaire}

\begin{center}
\begin{tabularx}{\linewidth}{|l|l|X|}
\hline
\rowcolor{popblueL} Mot anglais & Nature & Traduction française \\
\hline
findings & nom pluriel & résultats (d'une enquête) \\
\hline
nationwide & adjectif & à l'échelle nationale \\
\hline
sample & nom & échantillon \\
\hline
respondent & nom & personne interrogée \\
\hline
margin of error & locution nominale & marge d'erreur \\
\hline
turnout & nom & taux de participation \\
\hline
trend & nom & tendance \\
\hline
to take part & locution verbale & participer \\
\hline
\end{tabularx}
\end{center}

\textbf{Questions rapides} (répondez en anglais, en phrases complètes) :

\begin{enumerate}
\item Who conducted the survey, and where?
\item How many students were interviewed?
\item What percentage of respondents read for pleasure every week?
\item Why must the figures be read carefully?
\item What trend did the researchers notice?
\end{enumerate}

\begin{corrige}[Questions rapides]
\begin{enumerate}
\item \eng{A team of researchers from the University of Buea conducted the survey in secondary schools across Cameroon.}
\item \eng{Two thousand students were interviewed.}
\item \eng{Sixty per cent of respondents read for pleasure at least once a week.}
\item \eng{The figures must be read carefully because the margin of error is three per cent.}
\item \eng{They noticed that the older the students are, the more they read news articles rather than novels.}
\end{enumerate}
\end{corrige}

\courssub{Champ 1 — The actors (les acteurs)}

\begin{definition}[The people behind a poll]
Un \cle{sondage} met en relation plusieurs personnes : celui qui le conçoit, celui qui pose les questions et celui qui y répond.
\end{definition}

\begin{center}
\begin{tabularx}{\linewidth}{|X|X|X|}
\hline
\rowcolor{popblueL} English & Français & Exemple en contexte \\
\hline
pollster & sondeur (organisateur) & The pollster designed the questions. \\
\hline
researcher & chercheur & The researcher analysed the data. \\
\hline
interviewer & enquêteur (qui pose les questions) & The interviewer visited fifty homes. \\
\hline
respondent & personne interrogée & Each respondent gave honest answers. \\
\hline
\end{tabularx}
\end{center}

\begin{exemplebox}[Observez et traduisez]
\eng{The pollsters interviewed a sample of two thousand students.} « Les sondeurs ont interrogé un échantillon de deux mille élèves. »

\eng{Each respondent answered a short questionnaire.} « Chaque personne interrogée a répondu à un court questionnaire. »
\end{exemplebox}

\begin{attention}[Prononciation et vocabulaire]
\eng{researcher} se prononce « ri-seur-tcheur » (accent sur la deuxième syllabe). \eng{respondent} se prononce « ri-spon-dent ». Attention : en anglais, on ne dit jamais \eng{an interviewed} pour « une personne interrogée » ; le nom correct est \eng{a respondent} (ou \eng{an interviewee}).
\end{attention}

\courssub{Champ 2 — The tools (les outils)}

\begin{definition}[Poll, survey, census : ne pas confondre]
Un \cle{poll} est un sondage (souvent court, sur une opinion). Une \cle{survey} est une enquête (plus large, souvent statistique). Un \cle{census} est un recensement : on interroge \emph{toute} la population, pas un échantillon.
\end{definition}

\begin{center}
\begin{tabularx}{\linewidth}{|X|X|X|}
\hline
\rowcolor{popblueL} English & Français & Remarque \\
\hline
poll & sondage & souvent politique ou d'opinion \\
\hline
survey & enquête (statistique) & plus général qu'un poll \\
\hline
questionnaire & questionnaire & liste de questions écrites \\
\hline
form & formulaire & document à remplir \\
\hline
sample & échantillon & partie représentative de la population \\
\hline
census & recensement & compte toute la population \\
\hline
\end{tabularx}
\end{center}

\begin{exemplebox}[Exemples traduits]
\eng{They carried out a nationwide poll.} « Ils ont réalisé un sondage à l'échelle nationale. »

\eng{Only a sample of the population was questioned.} « Seul un échantillon de la population a été interrogé. »

\eng{A census counts every single person in the country.} « Un recensement compte chaque personne du pays. »
\end{exemplebox}

\begin{attention}[Faux amis dangereux]
« Une enquête » policière se dit \eng{an investigation}, jamais \eng{a survey} ! « Une enquête » statistique ou sociologique se dit \eng{a survey}. Autre piège classique : \eng{actually} signifie « en fait », et non « actuellement » ; pour « actuellement », dites \eng{currently} ou \eng{at present}.
\end{attention}

\courssub{Champ 3 — The actions (les actions)}

\begin{propriete}[Les verbes du sondage et leurs collocations]
En anglais, un sondage ne se « fait » pas (\eng{to make a survey} est un calque fautif du français « faire une enquête »). Les collocations correctes sont : \cle{to conduct a survey}, \cle{to carry out a survey}, \cle{to run a poll} (mener une enquête) ; \cle{to interview} (interroger) ; \cle{to answer} / \cle{to respond} (répondre) ; \cle{to take part in} (participer à) ; \cle{to publish findings} (publier les résultats).
\end{propriete}

\begin{exemplebox}[Exemples traduits]
\eng{The team conducted the survey over three months.} « L'équipe a mené l'enquête pendant trois mois. »

\eng{Almost every student agreed to take part.} « Presque tous les élèves ont accepté de participer. »

\eng{The findings will be published next month.} « Les résultats seront publiés le mois prochain. »
\end{exemplebox}

\begin{attention}[Prépositions piégées]
On dit \eng{to take part \textbf{in} a survey} (jamais « to take part to »). On dit \eng{to answer a question} (sans préposition) mais \eng{to respond \textbf{to} a question}. On dit \eng{to participate \textbf{in}}, pas « participate to ». De même : \eng{to listen \textbf{to}}, \eng{to wait \textbf{for}} — l'anglais choisit ses prépositions différemment du français.
\end{attention}

\courssub{Champ 4 — The results (les résultats)}

\begin{definition}[Turnout]
\cle{Turnout} = the percentage of people who vote or take part (taux de participation). Exemple : \eng{The turnout was excellent.} « Le taux de participation était excellent. »
\end{definition}

\begin{center}
\begin{tabularx}{\linewidth}{|X|X|X|}
\hline
\rowcolor{popblueL} English & Français & Exemple \\
\hline
results / findings & résultats & The findings surprised teachers. \\
\hline
percentage & pourcentage & A high percentage of students agreed. \\
\hline
majority & majorité & The majority prefer reading on phones. \\
\hline
minority & minorité & Only a minority read novels. \\
\hline
turnout & taux de participation & The turnout was excellent. \\
\hline
margin of error & marge d'erreur & The margin of error is three per cent. \\
\hline
trend & tendance & There is a clear trend towards online reading. \\
\hline
\end{tabularx}
\end{center}

\begin{exemplebox}[Exemples traduits]
\eng{Sixty per cent of respondents trust the media.} « Soixante pour cent des personnes interrogées font confiance aux médias. »

\eng{The majority of students read on their phones.} « La majorité des élèves lisent sur leur téléphone. »
\end{exemplebox}

\begin{attention}[Per cent, percentage et accords]
En anglais britannique, on écrit \eng{per cent} en deux mots (\eng{percent} en américain). Après \eng{per cent of + nom pluriel}, le verbe se met au pluriel : \eng{Sixty per cent of respondents \textbf{trust}...} Même logique avec \eng{the majority of} : \eng{the majority of students \textbf{prefer}} (pluriel). Ne confondez pas \eng{percentage} (le pourcentage, un chiffre : \eng{a high percentage}) et \eng{per cent} (après un nombre : \eng{fifty per cent}).
\end{attention}

\courssub{Carte mentale : les quatre champs lexicaux}

\begin{popfigure}
\begin{center}
\begin{tikzpicture}
\node[draw=popdark, fill=popblueL, rounded corners, font=\bfseries, align=center] (c) at (0,0) {The language\\of polls};
\node[draw=popblue, fill=popblue!20, rounded corners, align=center, font=\small] (a) at (-4.6,2.2) {\textbf{Actors}\\pollster\\researcher\\interviewer\\respondent};
\node[draw=poporange, fill=poporangeL, rounded corners, align=center, font=\small] (t) at (4.6,2.2) {\textbf{Tools}\\poll, survey\\questionnaire\\form, sample\\census};
\node[draw=popgreen, fill=popgreenL, rounded corners, align=center, font=\small] (ac) at (-4.6,-2.2) {\textbf{Actions}\\to conduct / carry out\\to interview\\to answer / respond\\to take part in\\to publish findings};
\node[draw=poppurple, fill=poppurpleL, rounded corners, align=center, font=\small] (r) at (4.6,-2.2) {\textbf{Results}\\results / findings\\percentage\\majority / minority\\turnout\\margin of error, trend};
\draw[-{Stealth}, thick, popblue] (c) -- (a);
\draw[-{Stealth}, thick, poporange] (c) -- (t);
\draw[-{Stealth}, thick, popgreen] (c) -- (ac);
\draw[-{Stealth}, thick, poppurple] (c) -- (r);
\end{tikzpicture}
\end{center}
\legende{Carte mentale du vocabulaire des sondages : quatre champs à mémoriser ensemble.}
\end{popfigure}

\courssub{Familles de mots : un mot, plusieurs formes}

\begin{propriete}[Dérivation : construire le vocabulaire]
L'anglais forme des familles de mots par dérivation. Les connaître multiplie votre vocabulaire sans effort :
\begin{itemize}
\item \eng{to respond} (verbe, répondre) $\rightarrow$ \eng{respondent} (nom de personne) $\rightarrow$ \eng{response} (nom de l'action, réponse) ;
\item \eng{to survey} (verbe, enquêter sur) $\rightarrow$ \eng{a survey} (nom, une enquête) ;
\item \eng{opinion} (nom) $\rightarrow$ \eng{opinion poll} (sondage d'opinion) ;
\item \eng{to interview} (verbe) $\rightarrow$ \eng{interviewer} (celui qui interroge) / \eng{interviewee} (celui qui est interrogé).
\end{itemize}
\end{propriete}

\begin{center}
\begin{tabular}{|l|l|l|}
\hline
\rowcolor{popblueL} Verbe & Personne & Action / chose \\
\hline
to respond & respondent & response \\
\hline
to interview & interviewer / interviewee & interview \\
\hline
to survey & surveyor & survey \\
\hline
to research & researcher & research \\
\hline
\end{tabular}
\end{center}

\begin{attention}[Le suffixe -er / -or / -ee]
Le suffixe \eng{-er} ou \eng{-or} désigne celui qui fait l'action (\eng{interviewer}, \eng{researcher}), tandis que \eng{-ee} désigne celui qui la subit (\eng{interviewee}, \eng{employee}). Ne confondez pas : le \eng{respondent} est celui qui \emph{répond}, pas celui qui pose les questions ! Remarque : \eng{surveyor} désigne le plus souvent un géomètre-expert ; pour « enquêteur », préférez \eng{interviewer} ou \eng{researcher}.
\end{attention}

\begin{savaistu}[Why is it called a “poll”?]
The word \eng{poll} originally meant “head” in old English — counting heads meant counting people, which is exactly what a poll does! Even today, journalists speak of \eng{going to the polls} for « aller aux urnes ». Le mot \eng{census}, lui, vient du latin : à Rome, le \emph{censor} comptait les citoyens pour les impôts et l'armée.
\end{savaistu}

\begin{exoresolu}[Complétez avec le mot du champ lexical qui convient]
Énoncé — Complétez chaque phrase avec : \emph{respondents, sample, turnout, margin of error, carry out, findings, majority, questionnaire}.
\begin{enumerate}
\item The NGO wants to \ldots\ a survey on water access in Douala.
\item Only a \ldots\ of 500 people was interviewed.
\item Each \ldots\ contained fifteen questions.
\item Most \ldots\ said they were satisfied with the project.
\item The \ldots\ of voters chose the first candidate.
\item The \ldots\ was low: only forty per cent of people took part.
\item The \ldots\ is two per cent, so the results are reliable.
\item The team published its \ldots\ in a national newspaper.
\end{enumerate}
\tcblower
Solution détaillée :
\begin{enumerate}
\item \eng{carry out} — collocation : \eng{to carry out a survey} (mener une enquête).
\item \eng{sample} — 500 personnes = un échantillon, pas toute la population.
\item \eng{questionnaire} — « each » exige le singulier ; c'est la liste de questions.
\item \eng{respondents} — les personnes interrogées (pluriel après \eng{most}).
\item \eng{majority} — « la majorité des électeurs » ; contraire : \eng{minority}.
\item \eng{turnout} — le taux de participation (fort ou faible : \eng{high/low turnout}).
\item \eng{margin of error} — la marge d'erreur, exprimée en pourcentage.
\item \eng{findings} — les résultats d'une recherche, presque toujours au pluriel.
\end{enumerate}
\end{exoresolu}

\begin{exercice}[À vous de jouer]
\begin{enumerate}
\item Traduisez en anglais : « Les sondeurs ont mené une enquête nationale et la majorité des personnes interrogées font confiance aux résultats. »
\item Trouvez l'intrus et justifiez : \eng{poll — survey — investigation — census}.
\item Formez le nom de personne à partir de : \eng{to research}, \eng{to interview} (celui qui pose les questions), \eng{to respond}.
\item En trois phrases, décrivez un sondage imaginaire mené dans votre lycée (utilisez au moins : \eng{sample, respondents, findings}).
\end{enumerate}
\end{exercice}

\begin{corrige}[Exercice « À vous de jouer »]
\begin{enumerate}
\item \eng{The pollsters carried out a nationwide survey, and the majority of respondents trust the results.}
\item L'intrus est \eng{investigation} : c'est une enquête \emph{policière} ; les trois autres désignent des comptages d'opinions ou de personnes.
\item \eng{researcher} ; \eng{interviewer} ; \eng{respondent}.
\item Modèle : \eng{Our English club conducted a survey on a sample of one hundred students. Most respondents said they prefer speaking activities. The findings will be published on the school noticeboard.}
\end{enumerate}
\end{corrige}

\begin{aretenir}[L'essentiel du vocabulaire des sondages]
Retenez les quatre champs : les \cle{acteurs} (\eng{pollster, researcher, interviewer, respondent}), les \cle{outils} (\eng{poll, survey, questionnaire, form, sample, census}), les \cle{actions} (\eng{to conduct / carry out a survey, to interview, to take part in, to publish findings}) et les \cle{résultats} (\eng{findings, percentage, majority, turnout, margin of error, trend}). Méfiez-vous des faux amis : \eng{a survey} n'est pas une enquête policière, et \eng{actually} ne veut pas dire « actuellement ». Ce vocabulaire vous servira immédiatement dans la section suivante, consacrée aux stratégies de lecture (\eng{skimming}, \eng{scanning}, déduction du sens des mots).
\end{aretenir}

\coursec{Méthode de lecture : skimming, scanning, guessing}

Dans la section précédente, nous avons enrichi notre vocabulaire du thème des sondages (\eng{poll}, \eng{survey}, \eng{respondent}, \eng{sample}, \eng{margin of error}…). Mais un bon vocabulaire ne suffit pas : à l'examen, vous devrez lire un texte inconnu, souvent long, en un temps limité. Il faut donc des \cle{stratégies de lecture} efficaces. Cette section vous apprend trois techniques complémentaires : le \cle{skimming} (lecture rapide pour l'idée générale), le \cle{scanning} (recherche d'une information précise) et la \cle{déduction du sens} d'un mot inconnu (\eng{guessing meaning from context}).

\courssub{Pourquoi des stratégies de lecture ?}

Observez ces deux comportements de lecteurs face au même article de presse :

\begin{exemplebox}[Deux lecteurs, deux méthodes]
\textbf{Reader A} reads the first word, then the second, translates every word into French in his head, stops at every unknown word, and after ten minutes he has only finished the first paragraph.\\
\textbf{Reader B} reads the title, the first sentence of each paragraph and the conclusion in two minutes. She knows what the text is about. Then she looks for the three pieces of information the questions ask for. She finishes in six minutes and answers correctly.
\end{exemplebox}

Le lecteur B n'est pas plus fort en anglais : il utilise simplement les bonnes stratégies. Voici le texte qui servira de support à toute la section. Lisez-le une première fois sans dictionnaire.

\begin{textebox}[Why Polls Matter — support text]
In many countries, opinion polls have become part of everyday political life. Before an election, pollsters interview a sample of voters and try to predict the result. In Cameroon, as in many African nations, surveys are also used to understand what citizens think about education, health and security.

A poll is only as good as its sample. If pollsters only ask rich people in big cities, the findings will not represent the whole population. That is why serious organisations choose respondents carefully, from different regions, ages and social backgrounds.

Voter turnout is another subject that surveys study closely. In one district, voter turnout was low: only 30\% voted. Researchers wanted to know why. Their survey showed that many young people did not believe their vote could change anything. “I answered the questionnaire honestly,” said one respondent, “because I want the government to know that we are disappointed.”

Critics sometimes say that polls influence voters instead of simply measuring opinions. When people read that a candidate is leading by twenty points, some may decide not to vote at all. Others may vote for the “winner” just to be on the winning side.

Despite these criticisms, polls remain useful tools. They give a voice to ordinary citizens between elections. The findings were published in the newspaper last week, and they have already started a national debate about youth participation. Whether we like them or not, surveys are now part of modern citizenship — and learning to read them critically is a skill every citizen needs.
\end{textebox}

\textbf{Glossaire :}

\begin{center}
\begin{tabularx}{\linewidth}{|l|l|X|}
\hline
\rowcolor{popblueL} \textbf{Mot anglais} & \textbf{Nature} & \textbf{Traduction} \\
\hline
pollster & nom & sondeur (personne qui réalise des sondages) \\
\hline
sample & nom & échantillon \\
\hline
findings & nom (pluriel) & résultats, conclusions (d'une enquête) \\
\hline
respondent & nom & personne interrogée, répondant \\
\hline
turnout & nom & taux de participation \\
\hline
questionnaire & nom & questionnaire \\
\hline
to lead by twenty points & expression & devancer de vingt points \\
\hline
background & nom & milieu, origine (sociale) \\
\hline
\end{tabularx}
\end{center}

\courssub{Skimming : lire vite pour l'idée générale}

\begin{definition}[Skimming]
Le \cle{skimming} (littéralement « écrémer ») consiste à lire un texte très rapidement — en deux ou trois minutes — pour en saisir l'\cle{idée générale} (\eng{the gist}), sans chercher à comprendre chaque mot.
\end{definition}

\begin{methode}[Comment faire du skimming en 2 minutes]
\begin{enumerate}
\item Lis le \textbf{titre} et l'éventuel sous-titre : ils annoncent le sujet.
\item Lis la \textbf{première phrase de chaque paragraphe} : en anglais journalistique, elle contient presque toujours l'idée principale du paragraphe (\eng{topic sentence}).
\item Lis la \textbf{dernière phrase} du texte (la conclusion).
\item Ignore tous les mots inconnus, les exemples et les détails chiffrés à ce stade.
\item Résume ensuite le texte en \textbf{une seule phrase}, en anglais si possible.
\end{enumerate}
\end{methode}

Appliquons au texte support. Les premières phrases des cinq paragraphes sont :

\begin{exemplebox}[Topic sentences du texte]
\eng{In many countries, opinion polls have become part of everyday political life.}\\
\eng{A poll is only as good as its sample.}\\
\eng{Voter turnout is another subject that surveys study closely.}\\
\eng{Critics sometimes say that polls influence voters instead of simply measuring opinions.}\\
\eng{Despite these criticisms, polls remain useful tools.}
\end{exemplebox}

À partir de ces cinq phrases seulement, on peut déjà produire un résumé global :

\begin{aretenir}[Idée générale du texte en une phrase]
\eng{The text explains how polls work, why the sample matters, what they reveal about voter turnout, why some people criticise them, and why they are still useful for citizens.}\\
« Le texte explique comment fonctionnent les sondages, pourquoi l'échantillon est important, ce qu'ils révèlent sur la participation électorale, pourquoi certains les critiquent et pourquoi ils restent utiles aux citoyens. »
\end{aretenir}

\courssub{Scanning : chercher une information précise}

\begin{definition}[Scanning]
Le \cle{scanning} consiste à parcourir le texte des yeux \textbf{sans le lire}, à la recherche d'un \cle{mot-clé} : un chiffre, un nom propre, une date, un pourcentage, un mot du thème. On s'arrête de lire en détail uniquement autour du mot-clé trouvé.
\end{definition}

\begin{methode}[Comment faire du scanning]
\begin{enumerate}
\item Lis d'abord la \textbf{question} et souligne le mot-clé à chercher (ex. : \eng{How many people voted?} → chercher un chiffre ou un \%).
\item Laisse tes yeux glisser sur le texte en diagonale, sans lire les phrases.
\item Quand tu repères le mot-clé (\eng{30\%}), lis \textbf{la phrase entière} qui l'entoure.
\item Formule ta réponse avec les mots de la question, sans recopier tout le paragraphe.
\end{enumerate}
\end{methode}

\begin{exemplebox}[Scanning appliqué]
Question : \eng{What percentage of people voted in the district?}\\
Les yeux cherchent un signe \% → on trouve : \eng{Voter turnout was low: only 30\% voted.}\\
Réponse courte : \eng{Only 30\% of the people voted.} « Seulement 30 \% des gens ont voté. »
\end{exemplebox}

\begin{exoresolu}[Entraînement chronométré : 3 questions de scanning en 3 minutes]
Énoncé — Répondez en 3 minutes maximum, sans relire tout le texte :
\begin{enumerate}
\item \eng{Who said: “I answered the questionnaire honestly”?}
\item \eng{By how many points is the candidate leading, according to the critics' example?}
\item \eng{Where were the findings published?}
\end{enumerate}
\tcblower
Solution détaillée :
\begin{enumerate}
\item Mot-clé : la citation entre guillemets. On la repère visuellement grâce aux guillemets. Réponse : \eng{A respondent (said it).} « Un répondant. »
\item Mot-clé : un nombre. On repère \eng{twenty points} dans le quatrième paragraphe. Réponse : \eng{The candidate is leading by twenty points.} « Le candidat devance de vingt points. »
\item Mot-clé : un lieu / un média. On repère \eng{newspaper} dans le dernier paragraphe. Réponse : \eng{The findings were published in the newspaper.} « Les résultats ont été publiés dans le journal. »
\end{enumerate}
\end{exoresolu}

\courssub{Guessing : déduire le sens d'un mot inconnu}

Même avec un excellent vocabulaire, vous rencontrerez toujours des mots inconnus. Trois outils permettent d'en deviner le sens : le \cle{contexte}, la \cle{famille de mots} et les \cle{préfixes et suffixes}.

\textbf{1. Utiliser le contexte.} Les mots voisins expliquent souvent le mot inconnu :

\begin{exemplebox}[Déduire « turnout » grâce au contexte]
\eng{Voter turnout was low: only 30\% voted.}\\
Les deux points annoncent une explication : « seulement 30 \% ont voté ». Donc \eng{turnout} désigne le \textbf{nombre de personnes qui votent}, c'est-à-dire le \textbf{taux de participation}. Vérification : « La participation électorale était faible : seulement 30 \% ont voté. » La phrase garde un sens parfait.
\end{exemplebox}

\textbf{2. Utiliser la famille de mots et les suffixes.} Un affixe révèle la nature et le sens du mot :

\begin{center}
\begin{tabularx}{\linewidth}{|l|l|X|}
\hline
\rowcolor{popblueL} \textbf{Affixe} & \textbf{Sens / fonction} & \textbf{Exemples} \\
\hline
-er / -or & personne qui fait l'action & \eng{teach → teacher} (enseignant), \eng{conduct → conductor} (conducteur), \eng{poll → pollster} (sondeur) \\
\hline
-ster & personne qui fait (variante de -er) & \eng{poll → pollster}, \eng{trick → trickster}, \eng{team → teamster} \\
\hline
-tion / -sion & nom d'action & \eng{participate → participation} (participation), \eng{decide → decision} (décision) \\
\hline
re- & de nouveau, à nouveau & \eng{count → recount} (recompter), \eng{elect → re-elect} (réélire) \\
\hline
-ist & spécialiste, adepte & \eng{journal → journalist} (journaliste), \eng{science → scientist} (scientifique) \\
\hline
-less & sans, privé de & \eng{hope → hopeless} (sans espoir) \\
\hline
un- & négation & \eng{fair → unfair} (injuste), \eng{likely → unlikely} (peu probable) \\
\hline
\end{tabularx}
\end{center}

\begin{exemplebox}[Exemple traité pas à pas : « pollster »]
\eng{Before an election, pollsters interview a sample of voters.}
\begin{enumerate}
\item \textbf{Nature} : le mot est sujet du verbe \eng{interview} et a un pluriel en -s → c'est un \textbf{nom de personne}.
\item \textbf{Famille} : \eng{poll} (sondage) + suffixe \eng{-ster} (variante de \eng{-er} : celui qui fait).
\item \textbf{Hypothèse} : « personne qui fait des sondages » = \textbf{un sondeur}.
\item \textbf{Vérification} : « Avant une élection, les sondeurs interrogent un échantillon d'électeurs. » Le sens est cohérent.
\end{enumerate}
\end{exemplebox}

\begin{exemplebox}[Autre exemple traduit]
\eng{The findings were published in the newspaper.}\\
\eng{findings} = \eng{results} → « Les résultats ont été publiés dans le journal. »\\
Indice : le suffixe \eng{-ing} forme ici un nom ; le contexte (\eng{published in the newspaper}) montre qu'il s'agit des conclusions d'une enquête.
\end{exemplebox}

\begin{methode}[Comment deviner le sens d'un mot inconnu]
\begin{enumerate}
\item \textbf{Identifie sa nature} (nom, verbe, adjectif, adverbe) grâce à sa place dans la phrase et à ses suffixes.
\item \textbf{Cherche sa famille} : connais-tu un mot de la même racine (\eng{poll} dans \eng{pollster}) ?
\item \textbf{Exploite le contexte} : les mots voisins, les deux-points, les exemples, les oppositions (\eng{but}, \eng{despite}).
\item \textbf{Remplace le mot par ton hypothèse} et vérifie que la phrase garde un sens.
\item Si tout échoue et que le mot est indispensable, consulte le glossaire ou le dictionnaire — en dernier recours seulement.
\end{enumerate}
\end{methode}

\begin{popfigure}
\begin{tikzpicture}[
node distance=0.9cm and 1.6cm,
q/.style={diamond, draw=popblue, fill=popblueL, align=center, inner sep=1.5pt, aspect=2, font=\small},
act/.style={rectangle, rounded corners, draw=poppurple, fill=poppurpleL, align=center, font=\small, inner sep=4pt},
fleche/.style={-{Stealth[length=2.5mm]}, thick, popdark}]
\node[q] (start) {Mot inconnu ?};
\node[q, below left=1.1cm and 0.2cm of start, align=center] (fam){Famille\\connue ?};
\node[act, right=1.8cm of fam, align=center] (ded){Déduire par la racine\\et le suffixe};
\node[q, below=1.1cm of fam, align=center] (ctx){Contexte\\suffisant ?};
\node[act, right=1.8cm of ctx, align=center] (hyp){Faire une hypothèse\\et vérifier le sens};
\node[act, below=1.1cm of ctx, align=center] (dic){Glossaire ou\\dictionnaire};
\draw[fleche] (start) -- (fam);
\draw[fleche] (fam) -- node[above, font=\small]{oui} (ded);
\draw[fleche] (fam) -- node[left, font=\small]{non} (ctx);
\draw[fleche] (ctx) -- node[above, font=\small]{oui} (hyp);
\draw[fleche] (ctx) -- node[left, font=\small]{non} (dic);
\end{tikzpicture}
\legende{Organigramme de décision face à un mot inconnu pendant la lecture.}
\end{popfigure}

\courssub{Comparaison des trois stratégies}

\begin{center}
\begin{tabularx}{\linewidth}{|l|X|X|X|}
\hline
\rowcolor{popblueL} \textbf{Stratégie} & \textbf{But} & \textbf{Comment} & \textbf{Quand l'utiliser} \\
\hline
\cle{Skimming} & Idée générale (\eng{gist}) & Titre, premières phrases, conclusion & Avant de répondre aux questions globales \\
\hline
\cle{Scanning} & Information précise & Chercher un mot-clé, un chiffre, un nom & Pour les questions de détail \\
\hline
\cle{Guessing} & Sens d'un mot inconnu & Nature, famille, contexte, vérification & Quand un mot bloque la compréhension \\
\hline
\end{tabularx}
\end{center}

\begin{attention}[Pièges fréquents des lecteurs francophones]
\begin{itemize}
\item \textbf{Ne traduis pas mot à mot} pendant la lecture : cela ralentit énormément et fait perdre le sens global. Lis par \cle{groupes de mots} (\eng{chunks}) : \eng{opinion polls / have become / part of everyday political life}, et non \eng{opinion - polls - have - become - part - of…}.
\item Ne t'arrête pas sur chaque mot inconnu : la plupart ne sont pas nécessaires pour répondre aux questions.
\item Attention aux \textbf{faux amis} dans ce thème : \eng{a survey} est « une enquête / un sondage », pas « une survie » ; \eng{to demand} signifie « exiger », pas « demander » (qui se dit \eng{to ask}) ; \eng{actually} signifie « en fait », pas « actuellement » (\eng{currently}).
\item Ne recopie pas de longues phrases du texte dans tes réponses : reformule brièvement avec tes propres mots.
\end{itemize}
\end{attention}

\begin{exoresolu}[Application complète : deviner trois mots du texte]
Énoncé — Sans dictionnaire, devinez le sens des mots soulignés et justifiez votre raisonnement :
\begin{enumerate}
\item \eng{Researchers wanted to know why. Their survey showed that many young people did not believe their vote could change anything.} → \eng{researchers}
\item \eng{Despite these criticisms, polls remain useful tools.} → \eng{despite}
\item \eng{Some may vote for the “winner” just to be on the winning side.} → \eng{winner}
\end{enumerate}
\tcblower
Solution détaillée :
\begin{enumerate}
\item \eng{research} (recherche) + suffixe \eng{-er} (personne qui fait) → « chercheurs ». Vérification : « Les chercheurs voulaient savoir pourquoi. » Sens correct.
\item \eng{Despite} est suivi d'un nom et introduit une opposition avec la suite (\eng{polls remain useful}) → « malgré ». Traduction : « Malgré ces critiques, les sondages restent des outils utiles. »
\item \eng{win} (gagner) + \eng{-er} → « le gagnant ». Le contexte (\eng{the winning side}, « le camp gagnant ») confirme l'hypothèse.
\end{enumerate}
\end{exoresolu}

\begin{exercice}[À vous de jouer]
\begin{enumerate}
\item Faites du skimming sur le texte support en 2 minutes, puis écrivez un résumé d'une phrase commençant par \eng{The text is about…}.
\item Scanning (2 minutes) : \eng{What three topics do surveys study in Cameroon, according to paragraph 1?}
\item Devinez le sens de \eng{disappointed} dans : \eng{“I want the government to know that we are disappointed.”} Justifiez votre réponse en trois étapes (nature, contexte, vérification).
\end{enumerate}
\end{exercice}

\begin{corrige}[Exercice « À vous de jouer »]
\begin{enumerate}
\item Réponse possible : \eng{The text is about the role of opinion polls in modern citizenship, their limits and their usefulness.} « Le texte porte sur le rôle des sondages d'opinion dans la citoyenneté moderne, leurs limites et leur utilité. »
\item On scanne le premier paragraphe à la recherche d'une énumération : \eng{Surveys are used to understand what citizens think about education, health and security.} « l'éducation, la santé et la sécurité ».
\item \textbf{Nature} : adjectif (après \eng{are}). \textbf{Contexte} : le répondant dit que les jeunes ne croient pas que leur vote peut changer les choses → sentiment négatif. \textbf{Hypothèse} : « déçus ». \textbf{Vérification} : « Je veux que le gouvernement sache que nous sommes déçus. » La phrase garde un sens cohérent.
\end{enumerate}
\end{corrige}

\begin{aretenir}[L'essentiel de la section]
\begin{itemize}
\item \cle{Skimming} : titre + premières phrases + conclusion → l'idée générale en 2 minutes.
\item \cle{Scanning} : repérer un mot-clé (chiffre, nom, date) sans tout lire → l'information précise.
\item \cle{Guessing} : nature du mot + famille (préfixes/suffixes) + contexte → hypothèse vérifiée dans la phrase.
\item On lit par \cle{groupes de mots}, jamais mot à mot, et le dictionnaire reste le dernier recours.
\end{itemize}
\end{aretenir}

\coursec{Production guidée : summary and opinion}

Après avoir maîtrisé les stratégies de lecture rapide (\emph{skimming}), de recherche ciblée (\emph{scanning}) et de déduction du sens (\emph{guessing}), vous êtes prêt à passer à l'écrit. L'épreuve de \emph{Reading} au Baccalauréat A se termine systématiquement par deux exercices de production écrite : un \textbf{résumé} (\emph{summary}) du texte lu et une \textbf{opinion personnelle} (\emph{personal opinion}) argumentée. Cette section vous guide pas à pas pour réussir ces deux exercices incontournables.

\courssub{Étape 1 — Résumer : l'art de la synthèse (Writing a Summary)}

Résumer, ce n'est pas recopier. C'est \textbf{restituer l'essentiel} (l'idée générale + les arguments principaux) \textbf{avec ses propres mots} (\emph{in your own words}) en un nombre de phrases imposé (ici 3 à 4 phrases). Le résumé doit être objectif : pas de « je », pas de jugement, juste ce que \emph{le texte dit}.

\begin{methode}[Comment rédiger un résumé efficace en 4 étapes]
\begin{enumerate}
\item \textbf{Relisez le texte une dernière fois} en surlignant \emph{qui} (sujet), \emph{quoi} (action/thème principal) et \emph{pourquoi/comment} (arguments clés). Ignorez les exemples secondaires, les chiffres précis et les anecdotes.
\item \textbf{Notez 3 ou 4 mots-clés} en anglais par paragraphe principal (ex. \eng{polls}, \eng{represent citizens}, \eng{influence policy}, \eng{example: election}).
\item \textbf{Redécrivez les idées en phrases complètes} en changeant la structure syntaxique et le vocabulaire (synonymes, passage de l'actif au passif, nominalisation). Visez 3 phrases maximum pour l'idée générale + 1 phrase pour les deux exemples regroupés.
\item \textbf{Comptez les mots et vérifiez} : accord sujet-verbe, connecteurs logiques (\eng{Firstly}, \eng{Moreover}, \eng{Finally}), ponctuation.
\end{enumerate}
\end{methode}

\begin{exemplebox}[Modèle de résumé basé sur le texte « Why Polls Matter » (Section 2)]
\eng{Opinion polls allow governments to measure public sentiment on key issues. They give a voice to ordinary citizens and help shape policies that reflect the majority's will. For instance, pre-election surveys guide parties in adjusting their campaigns, while satisfaction polls push local councils to improve public services like waste collection in Douala or road maintenance in Buea.}
\par\smallskip
\textbf{Traduction :} « Les sondages d'opinion permettent aux gouvernements de mesurer le sentiment public sur des enjeux majeurs. Ils donnent une voix aux citoyens ordinaires et aident à façonner des politiques qui reflètent la volonté de la majorité. Par exemple, les sondages pré-électoraux guident les partis dans l'ajustement de leurs campagnes, tandis que les sondages de satisfaction poussent les conseils municipaux à améliorer les services publics comme la collecte des déchets à Douala ou l'entretien des routes à Buea. »
\end{exemplebox}

\begin{attention}[Piège à éviter : le copier-coller (\emph{lifting})]
Ne recopiez \textbf{jamais} des phrases entières du texte. L'examinateur sanctionne le \emph{lifting} (copier-coller) par une note de 0 sur la compétence « Expression écrite ». \textbf{Paraphrasez} :
\begin{itemize}
    \item \eng{The text says: ``Polls are essential for democracy.''} $\rightarrow$ \textbf{Votre résumé :} \eng{The text highlights the crucial role of polls in a democracy.}
    \item \eng{``They help leaders make decisions.''} $\rightarrow$ \eng{They assist leaders in decision-making.}
\end{itemize}
\end{attention}

\begin{center}
\begin{tabularx}{\linewidth}{|X|X|X|}
\hline
\rowcolor{popblueL} \textbf{Original (Texte)} & \textbf{Mots-clés notés} & \textbf{Résumé (Vos mots)} \\
\hline
\eng{``Polls are a vital tool for democracy.''} & \eng{vital tool, democracy} & \eng{Polls serve as a vital democratic instrument.} \\
\hline
\eng{``They allow citizens to express views.''} & \eng{citizens express views} & \eng{They enable citizens to voice their opinions.} \\
\hline
\eng{``Example: the 2023 youth survey in Yaoundé.''} & \eng{example, youth survey, Yaoundé} & \eng{A recent youth survey in Yaoundé illustrates this.} \\
\hline
\end{tabularx}
\end{center}

\courssub{Étape 2 — Opinion : argumenter sa position (Giving your Opinion)}

Ici, le « je » est obligatoire (\eng{I think}, \eng{I believe}). La consigne type est : \eng{``Would you take part in a poll? Why or why not?''} (Participeriez-vous à un sondage ? Pourquoi ou pourquoi pas ?). Attendu : 4 à 5 phrases, au moins 2 connecteurs logiques, un argument \textbf{pour} et/ou \textbf{contre} nuancé.

\begin{definition}[Banque de formules : Exprimer son opinion]
Les formules ci-dessous sont vos « briques » pour construire un paragraphe d'opinion cohérent. Mémorisez-les.
\begin{center}
\begin{tabularx}{\linewidth}{|l|X|l|}
\hline
\rowcolor{popblueL} \textbf{Fonction} & \textbf{Formule anglaise} & \textbf{Équivalent français} \\
\hline
Introduire son avis & \eng{In my opinion,} / \eng{From my point of view,} & À mon avis / De mon point de vue \\
\hline
Affirmer une conviction & \eng{I strongly believe that...} / \eng{I am convinced that...} & Je crois fermement que... / Je suis convaincu que... \\
\hline
Donner une raison (cause) & \eng{... because / since / as...} & ... parce que / puisque / comme... \\
\hline
Nuancer (concession) & \eng{Although / Even though polls can be biased,} & Bien que / Même si les sondages puissent être biaisés, \\
\hline
Ajouter un argument & \eng{Furthermore,} / \eng{Moreover,} / \eng{In addition,} & De plus / En outre / En plus \\
\hline
Opposer / Contraster & \eng{On the one hand... on the other hand...} & D'un côté... de l'autre... \\
\hline
Conclure & \eng{That is why / Therefore / To sum up,} & C'est pourquoi / Par conséquent / En résumé \\
\hline
\end{tabularx}
\end{center}
\end{definition}

\begin{propriete}[Règle des connecteurs d'opinion : \eng{because} vs \eng{although}]
\begin{itemize}
    \item \textbf{\eng{because / since / as}} introduisent la \textbf{cause} (la raison de votre avis). \eng{I would participate because I want my voice heard.} (Je participerais \textbf{parce que} je veux qu'on m'entende.)
    \item \textbf{\eng{although / even though / though}} introduisent la \textbf{concession} (un fait qui va contre votre avis principal, mais ne le change pas). \eng{Although polls are sometimes inaccurate, I find them useful.} (\textbf{Bien que} les sondages soient parfois inexacts, je les trouve utiles.)
    \item \textbf{\eng{however / nevertheless}} s'utilisent en \textbf{début de phrase} (suivis d'une virgule) pour marquer l'opposition. \eng{Polls can be manipulated. However, they remain necessary.}
\end{itemize}
\end{propriete}

\begin{exemplebox}[Opinion modèle (4-5 phrases)]
\eng{In my opinion, I would definitely take part in a poll if I were asked. I believe that it is a civic duty to share one's views because decision-makers need reliable data to govern well. On the one hand, polls empower citizens; on the other hand, some people fear their data might be misused. Although this risk exists, I think the benefits of participating outweigh the drawbacks. That is why I would answer honestly.}
\par\smallskip
\textbf{Traduction :} « À mon avis, je participerais certainement à un sondage si on me le demandait. Je crois que c'est un devoir civique de partager son avis car les décideurs ont besoin de données fiables pour bien gouverner. D'un côté, les sondages donnent du pouvoir aux citoyens ; de l'autre, certains craignent que leurs données soient mal utilisées. Bien que ce risque existe, je pense que les avantages de la participation l'emportent sur les inconvénients. C'est pourquoi je répondrais honnêtement. »
\end{exemplebox}

\begin{attention}[Erreur fatale des francophones : \eng{``I am agree''} \textcolor{poporange}{$\times$}]
En anglais, \textbf{\eng{agree} est un VERBE}, pas un adjectif. On ne dit \textbf{jamais} \eng{``I am agree''}.
\begin{itemize}
    \item \textbf{Correct :} \eng{I agree.} / \eng{I don't agree.} / \eng{I disagree.}
    \item \textbf{Correct (avec adjectif) :} \eng{I am in agreement with you.} (Je suis d'accord avec toi.)
    \item \textbf{Correct (partiel) :} \eng{I partly agree.} / \eng{I agree to some extent.}
\end{itemize}
\cle{Agree = verbe} $\rightarrow$ \eng{Do you agree? Yes, I agree.}
\end{attention}

\courssub{Étape 3 — Relecture : la check-list finale (Proofreading Checklist)}

Avant de rendre votre copie, prenez 2 minutes pour vérifier ces 5 points (méthode \eng{COPS}) :

\begin{center}
\begin{tabular}{|c|p{10cm}|}
\hline
\rowcolor{poporangeL} \textbf{Lettre} & \textbf{Contrôle} \\
\hline
\textbf{C} & \textbf{Capitalization / Capitals} : Majuscules en début de phrase, pour \eng{I}, noms propres (\eng{Cameroon}, \eng{Yaoundé}, \eng{Monday}), nationalités (\eng{Cameroonian}). \\
\hline
\textbf{O} & \textbf{Organization / Order} : Ordre des mots \eng{Sujet + Verbe + Complément}. Connecteurs en place (\eng{However,} en début de phrase). Paragraphe unique ou deux paragraphes distincts (Résumé / Opinion). \\
\hline
\textbf{P} & \textbf{Punctuation} : Point final, virgule après connecteur introductif (\eng{In my opinion,}), apostrophe au génitif (\eng{citizens' voices}), pas d'espace avant le point. \\
\hline
\textbf{S} & \textbf{Spelling / Subject-Verb Agreement} : \eng{He thinks} (pas \eng{think}), \eng{The polls show} (pluriel), \eng{The government is} (singulier collectif). Vérifiez les irréguliers : \eng{make/made}, \eng{take/took}, \eng{give/given}. \\
\hline
\end{tabular}
\end{center}

\courssub{Complément Bac : le format « Summary + Opinion » à l'épreuve}

\begin{savaistu}[Le format « Summary + Opinion » au Baccalauréat A]
L'épreuve de \emph{Reading Comprehension} du Baccalauréat Série A se termine \textbf{systématiquement} par deux questions de production écrite valant ensemble \textbf{4 à 6 points} sur 20 :
\begin{enumerate}
    \item \textbf{Summary (2-3 pts) :} \eng{``Summarize the text in 3 or 4 sentences.''} ou \eng{``In about 40 words, summarize the main ideas of the passage.''}
    \item \textbf{Personal Opinion (2-3 pts) :} \eng{``Would you...? Why / Why not?''} ou \eng{``What is your opinion on...?''} (4-5 phrases attendues).
\end{enumerate}
\textbf{Conseil stratégique :} Traitez d'abord les questions de compréhension (faciles, points assurés), puis consacrez vos 15 dernières minutes à cette production. La qualité de l'anglais (grammaire, vocabulaire, connecteurs) prime sur la longueur.
\end{savaistu}

\courssub{Activité par binômes : Peer Review (Correction croisée)}

\begin{experience}[Échange et évaluation par les pairs]
\textbf{Consigne :} Échangez votre production (Résumé + Opinion) avec votre voisin. Lisez sa copie et cochez les critères ci-dessous. Soyez bienveillant mais exigeant : vous l'aidez à gagner des points au Bac !

\begin{center}
\begin{tabularx}{\linewidth}{|X|c|c|}
\hline
\rowcolor{popgreenL} \textbf{Critère de réussite} & \textbf{Oui} & \textbf{Non / À améliorer} \\
\hline
\textbf{RÉSUMÉ} (3-4 phrases) & & \\
\hline
L'idée générale du texte est présente (qui/quoi/pourquoi). & $\square$ & $\square$ \\
\hline
Les 2 exemples principaux sont mentionnés (regroupés en 1 phrase). & $\square$ & $\square$ \\
\hline
Aucune phrase n'est copiée telle quelle du texte (paraphrase). & $\square$ & $\square$ \\
\hline
\textbf{OPINION} (4-5 phrases) & & \\
\hline
La réponse directe à la question est claire (\eng{Yes, I would} / \eng{No, I wouldn't}). & $\square$ & $\square$ \\
\hline
Au moins \textbf{2 connecteurs logiques} sont utilisés (\eng{because, although, however, moreover...}). & $\square$ & $\square$ \\
\hline
L'argumentation est cohérente (raison + exemple personnel ou fait). & $\square$ & $\square$ \\
\hline
\textbf{LANGUE} & & \\
\hline
Accords Sujet-Verbe corrects (\eng{He thinks}, \eng{Polls are}). & $\square$ & $\square$ \\
\hline
Pas d'erreur \eng{``I am agree''}. & $\square$ & $\square$ \\
\hline
Ponctuation et majuscules respectées. & $\square$ & $\square$ \\
\hline
\end{tabularx}
\end{center}
\textbf{Feedback oral (2 min) :} Dites à votre partenaire : \eng{``Your summary is clear but you copied one sentence. Try to paraphrase.''} ou \eng{``Good use of 'although', but check subject-verb agreement in sentence 3.''}
\end{experience}

\courssub{Exercice résolu : De la lecture à l'écriture}

\begin{textebox}[Texte support pour l'exercice : « Youth and Online Surveys »]
\eng{Nowadays, young Cameroonians spend an average of four hours a day on their smartphones. A recent survey conducted by the National Institute of Statistics in Yaoundé and Douala reveals that 65\% of respondents aged 15-25 have participated in at least one online poll in the past year. These polls, often shared via WhatsApp groups or Facebook pages, cover topics ranging from favourite music artists to serious issues like youth unemployment and the quality of education.}

\eng{While some experts praise this digital engagement as a new form of citizenship, others warn against the reliability of such data. Professor Amadou Bello, a sociologist at the University of Buea, argues that ``online samples are rarely representative because they exclude young people without internet access in rural areas like the Far North.'' However, for many urban youths, clicking ``Vote'' feels like the easiest way to make their voice heard without joining a political party.}

\eng{\textbf{Glossaire :} \textit{respondent (n.) = personne interrogée ; representative (adj.) = représentatif ; rural areas (n.) = zones rurales ; digital engagement (n.) = engagement numérique.}}
\end{textebox}

\begin{exoresolu}[Rédiger un résumé (3-4 phrases) et une opinion (4-5 phrases)]
\textbf{Consigne :} \eng{1. Summarize the text in 3-4 sentences. 2. Would you participate in an online poll shared on WhatsApp? Why or why not? (4-5 sentences)}

\tcblower

\textbf{SOLUTION DÉTAILLÉE}

\textbf{Étape 1 : Analyse \& Mots-clés (Drafting)}
\begin{itemize}
    \item \textbf{Para 1 (Faits) :} Young Cameroonians, 4h/day phone, 65\% did online poll (WhatsApp/FB), topics: music to unemployment.
    \item \textbf{Para 2 (Débat) :} Experts divided: digital citizenship vs reliability (Prof Bello: not representative, excludes rural youth). Urban youth feel it's easy participation.
\end{itemize}

\textbf{Étape 2 : Rédaction du Résumé (Model Answer)}
\eng{A recent study shows that a majority of young Cameroonians in major cities participate in online polls shared on social media about various topics. However, experts disagree on the value of this digital engagement. While some see it as active citizenship, Professor Bello warns that these surveys are not representative because they exclude rural youth without internet access. Despite this limit, many urban youths view online voting as an easy way to express their opinions.}
\par\smallskip
\textit{(4 phrases. Idée générale + chiffre clé + débat d'experts + exemple/conclusion. Zero copier-coller. Vocabulaire varié : \eng{majority/65\%, major cities/Yaoundé-Douala, digital engagement/online participation, exclude/leave out}.)}

\textbf{Étape 3 : Rédaction de l'Opinion (Model Answer)}
\eng{In my opinion, I would participate in an online poll on WhatsApp if the topic interests me, like education or jobs. I believe it is a fast way to share my view because I do not have time for political meetings. Although I know these polls are not scientific and exclude my friends in villages without internet, I think they still show a trend. Moreover, it costs nothing to click a button. That is why I would vote, but I would not trust the results blindly.}
\par\smallskip
\textit{(5 phrases. Réponse directe \eng{I would}. Connecteurs : \eng{because, although, moreover, that is why}. Nuance : \eng{not scientific, exclude friends... but still show a trend}. Pas de \eng{I am agree}.)}

\textbf{Étape 4 : Auto-correction (COPS)}
\begin{itemize}
    \item \textbf{C} : \eng{Cameroonians, WhatsApp, Professor Bello, Far North} $\checkmark$
    \item \textbf{O} : \eng{SVO} respecté. \eng{Although..., I think...} $\checkmark$
    \item \textbf{P} : Virgules après \eng{In my opinion, Although..., Moreover,} $\checkmark$
    \item \textbf{S} : \eng{young Cameroonians... participate} (pluriel), \eng{Professor Bello... argues} (3e pers. sg), \eng{polls... are} $\checkmark$
\end{itemize}
\end{exoresolu}

\courssub{Entraînement individuel : À vous de jouer !}

\begin{exercice}[Production autonome — Simulacre Bac]
\eng{Read the text ``Youth and Online Surveys'' again.}
\eng{1. Summarize the text in \textbf{exactly 3 sentences}. (Focus: General idea + Expert debate + Youth reaction).}
\eng{2. \textbf{Personal opinion (5 sentences):} ``Online polls on social media are a waste of time because they are not scientific.'' \textbf{Do you agree or disagree?} Explain your view using \eng{although}, \eng{however} and \eng{because}.}
\end{exercice}

\begin{corrige}[Exercice Production autonome]
\textbf{1. Summary (Proposition) :}
\eng{Online polls on social media attract a large number of young urban Cameroonians who see them as an easy form of participation. Nevertheless, sociologists like Professor Bello criticize their lack of representativeness since rural populations are excluded. Consequently, while these surveys reflect urban trends, they cannot measure the opinion of the whole youth.}

\textbf{2. Opinion (Proposition) :}
\eng{I disagree with the idea that online polls are a waste of time. Although they are not scientific and lack representativeness, they allow young people to express themselves freely. Because traditional politics often ignores us, social media becomes our only platform. However, we should not consider these results as absolute truth. That is why I participate but remain critical.}
\par\smallskip
\textbf{Vérifiez votre copie avec la grille COPS et la grille d'échange par binômes (ci-dessus).}
\end{corrige}

\begin{popfigure}
\begin{tikzpicture}[
    node distance=1.5cm and 2cm,
    block/.style={draw=popblue, thick, rounded corners, fill=popblueL, text width=5.5cm, align=center, minimum height=3.5cm, font=\small},
    arrow/.style={-Stealth, thick, poporange, shorten >=2pt, shorten <=2pt},
    title/.style={font=\bfseries\large, text=popdark, node distance=0.5cm}
]
% Nœuds
\node[block, align=center] (summary){
    \textbf{SUMMARY}\\[0.3cm]
    \eng{(What the TEXT says)}\\[0.5cm]
    \begin{itemize}
        \item[$\bullet$] Objective / Neutral
        \item[$\bullet$] Main idea + Key examples
        \item[$\bullet$] \textbf{3-4 sentences max}
        \item[$\bullet$] \emph{Own words} (Paraphrase)
        \item[$\bullet$] No \eng{I}, No opinion
    \end{itemize}
};

\node[block, right=of summary, align=center] (opinion){
    \textbf{OPINION}\\[0.3cm]
    \eng{(What YOU think + Why)}\\[0.5cm]
    \begin{itemize}
        \item[$\bullet$] Subjective / Personal
        \item[$\bullet$] Direct answer (\eng{Yes/No})
        \item[$\bullet$] \textbf{4-5 sentences}
        \item[$\bullet$] Arguments + Examples
        \item[$\bullet$] Connectors: \eng{because, although, however, moreover}
    \end{itemize}
};

% Flèche
\draw[arrow] (summary.east) -- node[above, font=\bfseries\small, text=poppurple] {TRANSITION} (opinion.west);

% Titres au-dessus
\node[title, above=of summary] (t1) {\textcolor{popblue}{ÉTAPE 1}};
\node[title, above=of opinion] (t2) {\textcolor{poporange}{ÉTAPE 2}};

% Rappel Bac en bas
\node[below=1cm of summary.south, anchor=north, text width=12cm, align=center, font=\itshape\small, text=popmuted] (bac) {\textbf{RAPPEL BAC A :} Ces deux blocs correspondent aux deux dernières questions de l'épreuve de Reading Comprehension. Maîtrisez la structure : \eng{Summary} $\rightarrow$ \eng{Opinion}.};
\end{tikzpicture}
\legende{Schéma de la production guidée : du résumé objectif à l'opinion argumentée.}
\end{popfigure}

\newpage

\coursec{Activité d'intégration}

\courssub{Objectifs de l'activité}

À l'issue de cette activité d'intégration, l'élève sera capable de :

\begin{itemize}
\item mobiliser le \cle{lexique des sondages} (\eng{poll, survey, respondent, questionnaire, percentage, majority, opinion}) dans une tâche de communication authentique ;
\item rédiger un \cle{mini-questionnaire} de cinq questions claires et correctes sur la citoyenneté des jeunes ;
\item appliquer les \cle{stratégies de lecture} étudiées dans la leçon (skimming, scanning, déduction du sens) pour traiter rapidement un document authentique et les réponses recueillies ;
\item exprimer des \cle{résultats chiffrés} en anglais (\eng{70\% of respondents said that...}) et une \cle{opinion personnelle} argumentée ;
\item présenter oralement un court rapport de sondage devant la classe, avec une prononciation et une intonation soignées.
\end{itemize}

\courssub{Situation}

Le club d'anglais de ton lycée à Yaoundé prépare son magazine de fin d'année, \emph{The Young Citizen}. Le rédacteur en chef a lancé un appel à contributions : chaque classe de Terminale doit réaliser un mini-sondage sur le thème \emph{Youth and Citizenship} et publier un court article de résultats. Ton groupe a été choisi pour représenter ta classe. Voici la note d'information envoyée par le club :

\begin{textebox}[Note from the English Club — The Young Citizen magazine]
Dear students,

Our school magazine is preparing a special issue on youth and citizenship in Cameroon. We would like each class to carry out a short opinion poll and send us a report of about 120 words.

Your report should answer questions such as: Do young people know their rights? Would they vote at eighteen? Do they take part in community activities? Are they interested in politics?

Please include: the number of respondents, at least three results with percentages, and one sentence giving your group's opinion. The best reports will be published in the magazine and read aloud at the closing ceremony.

Deadline: Friday 15th. Good luck!

The editorial team
\end{textebox}

\begin{definition}[Glossaire utile]
\begin{itemize}
\item \eng{to carry out a poll} (verbe) : réaliser un sondage ;
\item \eng{a report} (nom) : un rapport, un compte rendu ;
\item \eng{a respondent} (nom) : une personne interrogée, un répondant ;
\item \eng{a deadline} (nom) : une date limite ;
\item \eng{to read aloud} (verbe) : lire à haute voix ;
\item \eng{an editorial team} (nom) : une équipe de rédaction.
\end{itemize}
\end{definition}

\courssub{Consignes}

Travail par groupes de quatre. L'activité se déroule en cinq étapes numérotées, de la conception du questionnaire à la présentation orale.

\begin{enumerate}
\item \textbf{Compréhension de la note (5 min).} Lis la note du club en appliquant le \cle{skimming} (lecture rapide pour saisir l'idée générale), puis le \cle{scanning} (recherche d'informations précises). Réponds mentalement aux trois questions : Who wrote the note? What must the report contain? What is the deadline?
\item \textbf{Rédaction du questionnaire (10 min).} Avec ton groupe, écris un mini-questionnaire de \cle{cinq questions} sur la citoyenneté des jeunes. Privilégie les questions fermées (\eng{Yes/No}) ou à choix simples pour faciliter le dépouillement. Exemples : \eng{Do you know your rights as a citizen? Would you vote at eighteen? Have you ever taken part in a community clean-up?}
\item \textbf{Passation et dépouillement (10 min).} Interrogez un autre groupe (vos \eng{respondents}). Notez les réponses, comptez-les et calculez les \cle{pourcentages} (par exemple, 6 répondants sur 8 donnent 75 \%).
\item \textbf{Rédaction du rapport (15 min).} Rédigez un paragraphe d'environ 120 mots : présentez le sondage (nombre de répondants), donnez au moins \cle{trois résultats chiffrés} avec des structures variées, et terminez par \cle{une phrase d'opinion} introduite par \eng{In our opinion...} ou \eng{We believe that...}.
\item \textbf{Présentation orale et compilation (10 min).} Un porte-parole par groupe lit le rapport à la classe. Le professeur compile tous les résultats au tableau sous forme de \eng{class poll} : la classe compare alors les tendances et commente les différences entre groupes.
\end{enumerate}

\begin{attention}[Pièges à éviter]
\begin{itemize}
\item Ne dis pas \eng{the 70\% of respondents} : pas d'article défini devant un pourcentage général. Écris \eng{70\% of respondents said that...}.
\item Attention à la concordance des temps après \eng{said that} : \eng{They said that they \emph{would} vote at eighteen} (et pas \emph{will}).
\item \eng{to make a survey} est un faux calque du français « faire un sondage » ; l'expression correcte est \eng{to carry out} ou \eng{to conduct a survey}.
\item Ne confonds pas \eng{opinion poll} (sondage d'opinion) et \eng{election} (élection) : un sondage ne fait pas élire, il mesure une opinion.
\item Prononce \eng{survey} « seur-veï » (nom, accent sur la première syllabe) et \eng{percentage} « per-sen-tidj ».
\end{itemize}
\end{attention}

\courssub{Ressources à mobiliser}

\begin{propriete}[Exprimer les résultats d'un sondage]
\begin{center}
\begin{tabularx}{\linewidth}{|X|X|}\hline
\rowcolor{popblueL} \textbf{Structure} & \textbf{Exemple} \\ \hline
\eng{X\% of respondents + verbe} & \eng{70\% of respondents know their rights.} \\ \hline
\eng{The majority of... / Most...} & \eng{Most students would vote at eighteen.} \\ \hline
\eng{A minority of... / Few...} & \eng{Few respondents take part in community work.} \\ \hline
\eng{... said that + proposition} & \eng{Half of them said that they trusted the courts.} \\ \hline
\eng{According to our survey,...} & \eng{According to our survey, young people are politically aware.} \\ \hline
\eng{X out of Y + nom} & \eng{Six out of eight students read the news.} \\ \hline
\end{tabularx}
\end{center}
\end{propriete}

\begin{aretenir}[Exprimer son opinion en conclusion]
\begin{itemize}
\item \eng{In our opinion, young citizens should learn about their rights at school.} « À notre avis, les jeunes citoyens devraient apprendre leurs droits à l'école. »
\item \eng{We believe that voting is a duty, not only a right.} « Nous croyons que voter est un devoir, pas seulement un droit. »
\item \eng{These results show that... / This proves that...} « Ces résultats montrent que... / Cela prouve que... »
\end{itemize}
\end{aretenir}

Rappelle-toi aussi des stratégies de lecture de la leçon : le \cle{skimming} pour saisir le sens global d'un document, le \cle{scanning} pour repérer rapidement les chiffres et les dates, et la \cle{déduction du sens} des mots inconnus grâce au contexte. Ces trois stratégies te seront indispensables à l'épreuve de compréhension de l'écrit du Baccalauréat.

\courssub{Proposition de production et corrigé commenté}

\begin{exoresolu}[Modèle de rapport de sondage]
Rédige le rapport final de ton groupe pour le magazine \emph{The Young Citizen} (environ 120 mots).
\tcblower
\textbf{Production modèle :}

\eng{Last week, our group carried out a short poll on youth and citizenship. We interviewed eight students from another Terminale class in our school.}

\eng{Our survey shows that 75\% of respondents know their basic rights as citizens. However, only 25\% of them have ever taken part in a community activity, such as a clean-up campaign in their neighbourhood. The majority of respondents, six out of eight, said that they would vote at eighteen if elections were organised. Surprisingly, a minority of students are interested in political debates on the radio.}

\eng{In our opinion, these results prove that young Cameroonians are ready to become active citizens, but schools should organise more civic activities to help them.}

\textbf{Commentaire des choix de langue :}
\begin{itemize}
\item \eng{carried out a short poll} : collocation correcte (et non \emph{made a poll}) ; le prétérit est utilisé car l'action est terminée et datée (\eng{last week}).
\item \eng{75\% of respondents know} : pas d'article devant le pourcentage ; le verbe est au présent car le résultat est présenté comme un fait général.
\item \eng{have ever taken part} : present perfect avec \eng{ever} pour une expérience non datée.
\item \eng{said that they would vote} : concordance des temps respectée (\eng{would} après un verbe de parole au prétérit) ; \eng{if elections were organised} : conditionnel de type 2 pour une hypothèse.
\item \eng{six out of eight} : alternative élégante au pourcentage pour varier le style.
\item \eng{However} et \eng{Surprisingly} : connecteurs qui organisent le texte et nuancent les résultats.
\item \eng{In our opinion... should organise} : opinion claire avec le modal \eng{should} pour formuler une recommandation, conformément à la consigne du magazine.
\end{itemize}
\end{exoresolu}

\courssub{Bilan de l'activité}

Cette tâche d'intégration a permis de réunir, dans une situation de communication réaliste, toutes les compétences travaillées dans la leçon :

\begin{itemize}
\item \textbf{Lecture} : la note du club a été traitée grâce au skimming et au scanning, comme un véritable document authentique ;
\item \textbf{Lexique} : le vocabulaire des sondages (\eng{poll, survey, respondent, percentage, majority, to carry out}) a été réemployé dans un contexte productif, ce qui favorise sa mémorisation ;
\item \textbf{Grammaire} : l'expression des pourcentages, la concordance des temps dans le discours rapporté et les modaux d'opinion ont été mobilisés à l'écrit comme à l'oral ;
\item \textbf{Expression écrite et orale} : chaque groupe a produit un rapport structuré (présentation, résultats, opinion) et l'a présenté devant la classe, entraînant la prise de parole en public ;
\item \textbf{Dimension citoyenne} : le thème \emph{Youth and Citizenship} a conduit les élèves à réfléchir, en anglais, à leurs droits et devoirs de jeunes citoyens camerounais.
\end{itemize}

La compilation finale au tableau (\eng{class poll}) offre une conclusion collective : en comparant les résultats des groupes, la classe découvre concrètement comment un sondage reflète — parfois avec surprise — l'opinion d'un échantillon, et pourquoi les enquêtes d'opinion jouent un rôle important dans la vie démocratique.

\newpage

\coursec{Méthodes et exercices résolus}

\courssub{Texte support : Why Polls Matter}

\begin{textebox}[Why Polls Matter — an article for young citizens]
Every election year, the same question comes back: can we trust opinion polls? Some citizens refuse to take part in them, either because they do not trust the organisations that conduct them, or because they fear that their answers will not remain anonymous. Yet, when they are well conducted, polls reflect public opinion and give ordinary people a voice.

How does a poll work? A polling organisation questions a sample — that is, a small group of people chosen to represent the whole population. For example, a sample of 1,000 respondents can give a good picture of what millions of citizens think. The respondents answer a questionnaire, often by phone or online, and the results are then published in newspapers or on television.

However, a poll is never a perfect photograph of reality. Every survey has a margin of error, usually around three per cent, which means that the results are only an estimate. Moreover, a biased question — a question that pushes people to answer in a particular way — can make the results unreliable. That is why serious pollsters use neutral questions and representative samples.

In a democracy, polls matter because they help leaders understand what citizens really want before taking decisions. But polls only work if people accept to take part in them and answer honestly. Refusing to answer is also a choice — and it is a choice that makes your voice disappear.
\end{textebox}

\begin{definition}[Glossaire du texte]
\begin{itemize}
\item \eng{a poll} (nom) : un sondage (d'opinion).
\item \eng{a survey} (nom) : une enquête, une étude.
\item \eng{a sample} (nom) : un échantillon.
\item \eng{a respondent} (nom) : un répondant, une personne interrogée.
\item \eng{a questionnaire} (nom) : un questionnaire.
\item \eng{a margin of error} (nom) : une marge d'erreur.
\item \eng{biased} (adjectif) : biaisé, partial, orienté.
\item \eng{a pollster} (nom) : un sondeur.
\item \eng{reliable} (adjectif) : fiable ; \eng{unreliable} : peu fiable.
\item \eng{to conduct / to carry out a poll} (verbe) : réaliser, mener un sondage.
\end{itemize}
\end{definition}

\courssub{Savoir-faire 1 : Lire et comprendre un texte sur les sondages}

\begin{methode}[Réussir la compréhension globale d'un texte sur les sondages]
\begin{enumerate}
\item \textbf{Lire le titre et la source d'abord.} Le titre \eng{Why Polls Matter} annonce un texte \emph{argumentatif} : l'auteur va défendre l'importance des sondages. Prédire le contenu avant de lire active la compréhension.
\item \textbf{Faire un skimming.} Lire rapidement le premier et le dernier paragraphe, puis la première phrase de chaque paragraphe (\eng{topic sentence}), pour dégager l'idée générale (\eng{main idea}) en une phrase.
\item \textbf{Repérer les mots du champ lexical du sondage} : \eng{poll, survey, respondents, sample, opinion, margin of error, to conduct, to take part in}. Ils confirment le thème et aident à suivre l'argumentation.
\item \textbf{Identifier la structure argumentative} : thèse (\eng{polls give citizens a voice}), arguments (\eng{for instance, moreover}), éventuelle concession (\eng{admittedly, however}), conclusion.
\item \textbf{Formuler la main idea en anglais}, sans recopier le texte : \eng{The text explains why opinion polls are useful in a democracy and how citizens can take part in them.}
\end{enumerate}
\end{methode}

\begin{exoresolu}[Compréhension globale — type Bac]
\textbf{Énoncé.} Read the text \emph{Why Polls Matter} and answer: (a) What type of text is it: narrative, descriptive or argumentative? (b) State the main idea of the text in one sentence of your own. (c) Who is the text probably written for?
\tcblower
\textbf{Raisonnement.} (a) Le texte défend une idée (les sondages sont utiles à la démocratie) avec des arguments et des connecteurs comme \eng{moreover}, \eng{however} : c'est un texte \emph{argumentatif}, pas un récit. (b) La main idea doit résumer thèse + arguments, sans détail ni copie. (c) Le ton informatif et le vocabulaire civique (\eng{citizens, democracy}) visent un public général, par exemple des lycéens ou des lecteurs de journaux.

\textbf{Réponse rédigée.}
\eng{(a) It is an argumentative text, because the writer defends the idea that polls are important and supports it with arguments and examples.}
\eng{(b) The main idea is that opinion polls give ordinary citizens a voice in a democracy, provided that people take part in them honestly and that the results are read carefully.}
\eng{(c) The text is probably written for the general public, especially young citizens and students who are learning about civic life.}

\textbf{Conclusion.} Une bonne réponse de compréhension globale est courte, en anglais correct, et toujours \emph{justifiée} par un élément du texte (\eng{because...}).
\end{exoresolu}

\courssub{Savoir-faire 2 : Répondre aux questions de compréhension (scanning)}

\begin{methode}[Localiser et justifier : la technique du scanning]
\begin{enumerate}
\item \textbf{Lire la question avant le texte} et souligner les mots-clés : noms propres, dates, chiffres, verbes d'action.
\item \textbf{Scanner le texte} : faire glisser les yeux pour repérer le mot-clé ou un synonyme (attention aux paraphrases : \eng{take part in} = \eng{participate in}).
\item \textbf{Recopier la phrase du texte qui contient la réponse}, puis reformuler en phrase complète qui reprend les termes de la question.
\item \textbf{Citer la preuve} quand la consigne dit \eng{justify / quote from the text} : donner la phrase exacte entre guillemets anglais “ … ”.
\item \textbf{Pour les questions True/False}, écrire \eng{True} ou \eng{False} \emph{plus} la justification tirée du texte ; une réponse sans justification perd la moitié des points.
\item \textbf{Pour les questions de référence} (\eng{What does “they” refer to?}), remonter dans la phrase ou la phrase précédente pour trouver le nom auquel le pronom renvoie.
\end{enumerate}
\end{methode}

\begin{exoresolu}[Questions de compréhension — type Bac]
\textbf{Énoncé.} Answer with complete sentences and justify from the text.
(a) According to the text, why do some citizens refuse to take part in polls?
(b) True or False? Justify: \emph{Poll results are always perfectly exact.}
(c) What does the word “they” refer to in the sentence: \emph{“When they are well conducted, polls reflect public opinion”}?
\tcblower
\textbf{Raisonnement.} (a) Scanner les mots-clés \eng{refuse / take part} : le premier paragraphe donne deux raisons, la méfiance (\eng{distrust}) et la peur que les réponses ne restent pas anonymes. (b) Le troisième paragraphe mentionne la \eng{margin of error} : les sondages comportent une marge d'erreur, donc l'affirmation est fausse. (c) Le pronom \eng{they} est sujet du verbe \eng{are conducted} ; le seul nom pluriel logique est \eng{polls}.

\textbf{Réponses rédigées.}
\eng{(a) Some citizens refuse to take part in polls because they do not trust the organisations that conduct them, or because they fear that their answers will not remain anonymous.}
\eng{(b) False. The text says that every survey has a “margin of error”, which means that the results are only an estimate of public opinion, not an exact picture.}
\eng{(c) “They” refers to “polls”.}

\textbf{Conclusion.} Chaque réponse reprend les termes de la question, est rédigée en phrase complète et s'appuie sur une preuve textuelle citée entre guillemets.
\end{exoresolu}

\courssub{Savoir-faire 3 : Déduire le sens des mots inconnus}

\begin{methode}[Guessing meaning from context]
\begin{enumerate}
\item \textbf{Identifier la nature du mot} (nom, verbe, adjectif) grâce à sa position : \eng{a survey} (article + nom), \eng{they conduct} (sujet + verbe).
\item \textbf{Exploiter les indices du contexte} : synonyme (\eng{poll or survey}), contraire introduit par \eng{but, unlike}, définition (\eng{a sample — that is, a small group...}), exemple (\eng{such as}).
\item \textbf{Analyser la morphologie} : préfixes (\eng{dis-}trust = manque de confiance ; \eng{un-}reliable = peu fiable), suffixes (\eng{-er} : \eng{interviewer} = celui qui interviewe ; \eng{-less} : \eng{meaningless} = sans signification).
\item \textbf{Proposer un équivalent anglais simple} (\eng{a sample = a small representative group}) avant de traduire en français (« un échantillon »).
\item \textbf{Vérifier} en replaçant le sens trouvé dans la phrase : elle doit rester logique et grammaticale.
\end{enumerate}
\end{methode}

\begin{exoresolu}[Vocabulary in context — type Bac]
\textbf{Énoncé.} Find in the text words or expressions meaning: (a) a small group of people chosen to represent the whole population; (b) to carry out, to organise (a poll); (c) the possible difference between the poll results and reality. Then explain in your own words what “a biased question” means.
\tcblower
\textbf{Raisonnement.} (a) Le contexte \eng{a sample — that is, a small group of people} donne une définition explicite : \eng{a sample}. (b) Le verbe associé à \eng{poll} dans le texte est \eng{to conduct} (\eng{well conducted}). (c) L'expression technique du troisième paragraphe est \eng{margin of error}. Pour \eng{biased}, morphologie : \eng{bias} (parti pris) + suffixe adjectival \eng{-ed} ; le contexte (\eng{a question that pushes people to answer in a particular way}) confirme le sens « orientée, partiale ».

\textbf{Réponses rédigées.}
\eng{(a) “a sample” (b) “to conduct” (c) “margin of error”.}
\eng{A biased question is a question that is not neutral: it influences the respondent and pushes him or her towards a particular answer, so the results of the poll become unreliable.}

\textbf{Conclusion.} Ne jamais traduire mot à mot : le sens se déduit du contexte, puis se reformule en anglais simple.
\end{exoresolu}

\courssub{Savoir-faire 4 : Maîtriser le lexique des sondages}

\begin{methode}[Employer correctement le vocabulaire des polls]
\begin{enumerate}
\item \textbf{Apprendre les collocations}, pas les mots isolés : \eng{to conduct / to carry out a poll}, \eng{to take part in a survey}, \eng{to answer a questionnaire}, \eng{a representative sample}, \eng{public opinion}, \eng{the margin of error}.
\item \textbf{Distinguer les paires proches} : \eng{poll} (sondage d'opinion, souvent politique) / \eng{survey} (enquête plus large) ; \eng{respondent} (personne qui répond) / \eng{interviewer} ou \eng{pollster} (personne qui pose les questions).
\item \textbf{Attention aux prépositions} : \eng{to take part \textbf{in} a survey}, \eng{to be interested \textbf{in} politics}, \eng{according \textbf{to} the poll}. En revanche, \eng{to answer a question} se construit \emph{sans} préposition (jamais \eng{to answer to a question}).
\item \textbf{Réutiliser le lexique dans des phrases personnelles} liées au contexte camerounais : \eng{Last month, students of our school took part in a survey about the use of mobile phones.}
\end{enumerate}
\end{methode}

\begin{center}
\begin{tabularx}{\linewidth}{|X|X|X|}
\hline
\rowcolor{popblueL} English & Français & Exemple \\
\hline
a poll / an opinion poll & un sondage (d'opinion) & \eng{The latest poll shows a change in public opinion.} \\
\hline
a survey & une enquête & \eng{Our class carried out a survey on reading habits.} \\
\hline
a respondent & un répondant & \eng{Each respondent answered ten questions.} \\
\hline
a sample & un échantillon & \eng{A sample of 1,000 people was interviewed.} \\
\hline
a questionnaire & un questionnaire & \eng{Please fill in this questionnaire.} \\
\hline
the margin of error & la marge d'erreur & \eng{The margin of error is three per cent.} \\
\hline
public opinion & l'opinion publique & \eng{Public opinion is divided on this issue.} \\
\hline
to conduct / to carry out a poll & mener un sondage & \eng{Journalists conducted a poll in Douala.} \\
\hline
to take part in a survey & participer à une enquête & \eng{Will you take part in our survey?} \\
\hline
a biased question & une question orientée & \eng{A biased question makes the results unreliable.} \\
\hline
\end{tabularx}
\end{center}

\begin{exoresolu}[Lexique — complétion et réemploi]
\textbf{Énoncé.} Complete with: \emph{sample, respondents, conducted, margin of error, take part in, opinion}.
(a) The poll was \ldots by a team of journalists in Yaoundé.
(b) Only 500 \ldots answered the questionnaire.
(c) According to the survey, public \ldots is in favour of the new law.
(d) Would you like to \ldots our survey on road safety?
(e) The results are not exact: there is a \ldots of three per cent.
\tcblower
\textbf{Raisonnement.} (a) Structure passive : \eng{was conducted} (participe passé de \eng{conduct}). (b) Après un nombre, nom pluriel désignant les personnes interrogées : \eng{respondents}. (c) Collocation figée : \eng{public opinion}. (d) Après \eng{would you like to}, base verbale : \eng{take part in} (avec la préposition \eng{in} obligatoire). (e) Expression technique : \eng{margin of error}.

\textbf{Réponses.}
\eng{(a) conducted (b) respondents (c) opinion (d) take part in (e) margin of error.}

\textbf{Conclusion.} La bonne réponse dépend de la grammaire (passif, pluriel, base verbale après \eng{to}) autant que du sens : toujours vérifier la forme du mot placé.
\end{exoresolu}

\courssub{Savoir-faire 5 : Résumer et donner son opinion}

\begin{methode}[Summary and personal opinion en 5 à 8 phrases]
\begin{enumerate}
\item \textbf{Le résumé (3-4 phrases)} : reprendre uniquement les idées principales, avec ses propres mots, à la troisième personne, sans exemple ni chiffre secondaire. Amorce type : \eng{The text deals with... / The writer argues that...}
\item \textbf{Les connecteurs d'enchaînement} : \eng{first, moreover, however, finally}.
\item \textbf{L'opinion (2-3 phrases)} : s'exprimer à la première personne avec des structures variées : \eng{In my opinion... / I believe that... / As far as I am concerned... / I agree with the writer because...}
\item \textbf{Justifier l'opinion} par un argument ou une expérience personnelle : \eng{For instance, in my school, a survey showed that...}
\item \textbf{Relire} : vérifier l'accord sujet-verbe (\eng{polls \emph{are}}, pas \eng{polls is}), le présent de vérité générale, et l'absence de copie mot à mot du texte.
\end{enumerate}
\end{methode}

\begin{exoresolu}[Production guidée — summary and opinion]
\textbf{Énoncé.} In 60–80 words, summarise the text \emph{Why Polls Matter} and give your opinion about taking part in surveys.
\tcblower
\textbf{Raisonnement.} Résumé : idée 1 (les sondages donnent la parole aux citoyens), idée 2 (ils aident les dirigeants à prendre des décisions), idée 3 (ils doivent être bien menés pour être fiables). Opinion : position claire + justification personnelle. Temps : présent simple (vérités générales). Compter les mots pour respecter la consigne.

\textbf{Réponse rédigée (environ 75 mots).}
\eng{The text explains that opinion polls are an important tool in a democracy because they allow ordinary citizens to express their views and help leaders make better decisions. However, polls are only reliable when the sample is representative and the questions are neutral. In my opinion, taking part in surveys is a civic duty: if people refuse to answer, the results do not reflect reality. That is why I always answer questionnaires honestly.}

\textbf{Conclusion.} Un bon paragraphe de production combine résumé objectif (présent, troisième personne) et opinion justifiée (\eng{in my opinion} + \eng{because / that is why}).
\end{exoresolu}

\begin{attention}[Les 5 erreurs qui coûtent des points au Bac]
\begin{enumerate}
\item \textbf{Faux ami :} \eng{to assist} ne veut pas dire « assister à » au sens de « participer » (il signifie « aider »). Pour « participer à un sondage », écrire \eng{to take part in a poll} ou \eng{to participate in a survey}, jamais \eng{to assist a poll}.
\item \textbf{Préposition oubliée :} on écrit \eng{to take part \textbf{in} a survey}, \eng{to be interested \textbf{in} politics}, \eng{according \textbf{to} the results}. Le calque du français (« participer un sondage ») produit \eng{take part a survey}, qui est faux.
\item \textbf{Accord sujet-verbe :} \eng{public opinion \emph{is}} (singulier), mais \eng{the results \emph{are}}, \eng{polls \emph{show}}. Attention au \eng{-s} de la troisième personne : \eng{a poll \emph{shows}}, pas \eng{a poll show}.
\item \textbf{Copie intégrale du texte :} dans un résumé ou une réponse « in your own words », recopier des phrases entières du texte est sanctionné. Reformuler avec des synonymes : \eng{citizens} → \eng{ordinary people}, \eng{important} → \eng{useful}.
\item \textbf{Orthographe du lexique :} \eng{survey} (pas \eng{survay}), \eng{opinion} (un seul \eng{p}), \eng{questionnaire} (\eng{qu-} + double \eng{n}), \eng{respondent} (pas \eng{respondant}, calque du français « répondant »).
\end{enumerate}
\end{attention}

\begin{exercice}[Pour s'entraîner — reading et lexique]
\begin{enumerate}
\item Read the text again and say whether these sentences are True or False. Justify with a quotation.
(a) All citizens agree to answer polls.
(b) A sample must represent the whole population.
(c) Polls help leaders take decisions.
\item Find in the text a synonym for: (a) to organise (a poll); (b) not reliable; (c) ordinary people.
\item Complete with your own words: \eng{In my opinion, young Cameroonians should take part in surveys because...}
\end{enumerate}
\end{exercice}

\begin{corrige}[Exercice : Pour s'entraîner]
\begin{enumerate}
\item (a) \eng{False. The text says that “some citizens refuse to take part in them”.} (b) \eng{True. A sample is “a small group of people chosen to represent the whole population”.} (c) \eng{True. “Polls matter because they help leaders understand what citizens really want before taking decisions.”}
\item (a) \eng{to conduct / to carry out} ; (b) \eng{unreliable} ; (c) \eng{ordinary people / citizens}.
\item Réponse libre, par exemple : \eng{In my opinion, young Cameroonians should take part in surveys because it is a simple way to express their views and to influence the decisions of their leaders.}
\end{enumerate}
\end{corrige}

\newpage

\coursec{Exercices}

\courssub{Je vérifie mes connaissances}

\begin{exercice}[QCM : le bon mot]
Choisis la bonne réponse pour compléter chaque phrase.
\begin{enumerate}
\item A \ldots\ is a survey in which people are asked about their opinions.\\
a) census \quad b) poll \quad c) investigation
\item A \ldots\ counts every single person in a country.\\
a) sample \quad b) census \quad c) questionnaire
\item The people who answer a survey are called \ldots\\
a) respondents \quad b) reporters \quad c) candidates
\item A \ldots\ is a small group of people chosen to represent a larger population.\\
a) margin \quad b) turnout \quad c) sample
\item The \ldots\ of error tells us how accurate the results of a poll are.\\
a) rate \quad b) margin \quad c) level
\end{enumerate}
\end{exercice}

\begin{exercice}[Vrai ou faux justifié]
D'après le texte \emph{Why Polls Matter} étudié en classe, dis si chaque affirmation est \textbf{true} ou \textbf{false}, puis cite la ligne du texte qui justifie ta réponse.
\begin{enumerate}
\item Polls only interest politicians.
\item A well-conducted poll can reflect the opinion of millions of people.
\item The size of the sample does not matter at all.
\item Polls can influence public debate before an election.
\item All polls published in the media are reliable.
\end{enumerate}
\end{exercice}

\begin{exercice}[Vocabulaire : mots croisés du thème]
Complète chaque phrase avec un mot du lexique des sondages : \emph{survey, question, results, opinion, majority, publish}.
\begin{enumerate}
\item The \ldots\ of the poll will be announced on television tonight.
\item Sixty per cent of the \ldots\ think the government should build more schools.
\item Our class conducted a \ldots\ on students' eating habits.
\item Journalists often \ldots\ poll results before elections.
\item In my \ldots, young people should vote at eighteen.
\end{enumerate}
\end{exercice}

\begin{exercice}[Conjugaison : complète au bon temps]
Mets le verbe entre parenthèses au temps qui convient (présent simple, present perfect, futur, passif).
\begin{enumerate}
\item Every year, thousands of Cameroonians \ldots\ (take) part in surveys.
\item The results of the poll \ldots\ (publish) in yesterday's newspaper.
\item So far, the institute \ldots\ (interview) two thousand respondents.
\item Next month, a new survey \ldots\ (carry out) in our region.
\item People usually \ldots\ (answer) polls more honestly when they remain anonymous.
\end{enumerate}
\end{exercice}

\courssub{Je m'entraîne}

\begin{exercice}[Traduction]
Traduis les phrases suivantes en anglais.
\begin{enumerate}
\item Les sondages permettent aux citoyens d'exprimer leur opinion.
\item L'institut a interrogé un échantillon de mille personnes à Douala.
\item La marge d'erreur de ce sondage est de trois pour cent.
\item Beaucoup de jeunes refusent de répondre aux enquêtes par téléphone.
\item Les résultats seront publiés avant les élections.
\end{enumerate}
\end{exercice}

\begin{exercice}[Transformation de phrases]
Transforme chaque phrase selon la consigne entre parenthèses, sans changer le sens.
\begin{enumerate}
\item The institute interviewed 500 voters in Yaoundé. (passif)
\item “Do you trust the results of this poll?” the journalist asked me. (discours rapporté)
\item The sample is too small. The results are not reliable. (relier avec \emph{because})
\item People participate in polls willingly. (mettre à la forme négative)
\item The survey was conducted last week. (poser la question sur \emph{last week})
\end{enumerate}
\end{exercice}

\begin{exercice}[Texte à trous]
Complète le texte suivant avec les mots de la liste : \emph{respondents, sample, conducted, margin, turnout, questionnaire, accurate, opinion}.

\begin{textebox}[A survey in Buea]
Last month, a research centre \ldots\ a survey on young people's reading habits in Buea. A \ldots\ of 800 students was selected from five secondary schools. Each student received a \ldots\ with twenty questions. Most \ldots\ answered honestly, so the results are considered \ldots. The \ldots\ of error was only two per cent. The survey showed that the \ldots\ of students prefer reading on their phones, and in their \ldots, libraries should offer more digital books. The organisers hope for a high \ldots\ when the survey is repeated next year.
\end{textebox}
\end{exercice}

\begin{exercice}[Appariements : mots et définitions]
Associe chaque mot (1--8) à sa définition (a--h).
\begin{enumerate}
\item respondent
\item turnout
\item margin of error
\item sample
\item census
\item questionnaire
\item to conduct
\item forecast
\end{enumerate}
\begin{enumerate}[label=\alph*)]
\item a prediction about what will happen, for example in an election
\item a person who answers the questions of a survey
\item an official count of all the people in a country
\item the number or percentage of people who take part in something
\item a small group chosen to represent a larger population
\item a written list of questions used in a survey
\item the degree to which the results of a poll may be wrong
\item to organise and carry out an activity
\end{enumerate}
\end{exercice}

\courssub{Je me prépare au Bac}

\begin{exercice}[Type Bac : compréhension écrite]
Lis le texte suivant, puis réponds aux questions.

\begin{textebox}[Poll fever in Nigeria]
As Nigeria's general election approaches, opinion polls have become a daily topic of conversation. Every week, research institutes publish new figures showing which candidate is leading. In Lagos, Africa's largest city, interviewers stand at bus stops and in markets, clipboard in hand, asking people about their voting intentions.

Not everyone trusts these polls. Mrs Adebayo, a trader in Balogun market, refuses to answer. “Last time, the polls said one candidate would win easily, and he lost,” she explains. “How can they ask two thousand people and claim to know what two hundred million think?”

Pollsters reply that a carefully selected sample can indeed represent a whole nation, provided the margin of error is clearly stated. They also admit that some voters change their minds at the last minute, or hide their real intentions from interviewers.

Despite the criticism, polls play a useful role. They push candidates to address the issues voters care about, such as jobs, security and the cost of food. They also encourage citizens to discuss politics openly. Whether accurate or not, polls have made this election the most talked-about in the country's recent history.
\end{textebox}

\textbf{Questions}
\begin{enumerate}
\item Where do interviewers in Lagos ask people about their voting intentions?
\item Why does Mrs Adebayo refuse to answer pollsters' questions?
\item What condition must be met for a sample to represent a whole nation? Answer \textbf{in your own words}.
\item Mention two reasons why poll results can be wrong, according to the text.
\item Give two positive effects of polls mentioned in the last paragraph.
\item Find in the text : (a) a synonym of “figures”; (b) the opposite of “hide”.
\end{enumerate}
\end{exercice}

\begin{exercice}[Type Bac : traduction]
Traduis en français le passage suivant, extrait du texte \emph{Poll fever in Nigeria} :

\begin{textebox}
“Pollsters reply that a carefully selected sample can indeed represent a whole nation, provided the margin of error is clearly stated. They also admit that some voters change their minds at the last minute, or hide their real intentions from interviewers.”
\end{textebox}
\end{exercice}

\begin{exercice}[Type Bac : expression écrite]
Your school wants to conduct a survey on students' use of social media. Write a short paragraph (60--80 words) explaining the purpose of the survey and giving your own opinion about it. Use at least three words from the poll vocabulary studied in this lesson (e.g. \emph{survey, respondents, sample, results, opinion}).
\end{exercice}

\newpage

\coursec{Corrigés des exercices}

\begin{corrige}[Exercice 1 — QCM : le bon mot]
\begin{enumerate}
\item \textbf{b) poll} — Un \eng{poll} est un sondage d'opinion ; \eng{census} = recensement, \eng{investigation} = enquête.
\item \textbf{b) census} — Un \eng{census} compte toute la population ; \eng{sample} = échantillon, \eng{questionnaire} = questionnaire.
\item \textbf{a) respondents} — Les personnes interrogées sont les \eng{respondents} (répondants) ; \eng{reporters} = journalistes, \eng{candidates} = candidats.
\item \textbf{c) sample} — Un \eng{sample} est un échantillon représentatif ; \eng{margin} = marge, \eng{turnout} = taux de participation.
\item \textbf{b) margin} — La \eng{margin of error} (marge d'erreur) indique la précision ; \eng{rate} = taux, \eng{level} = niveau.
\end{enumerate}
\end{corrige}

\begin{corrige}[Exercice 2 — Vrai ou faux justifié]
\begin{enumerate}
\item \textbf{False.} Le texte explique que les sondages intéressent aussi les citoyens, les médias et les chercheurs (l. 2‑4).
\item \textbf{True.} Le texte affirme qu'un sondage bien mené peut refléter l'opinion de millions de personnes grâce à un échantillon représentatif (l. 6‑8).
\item \textbf{False.} La taille de l'échantillon est cruciale : un échantillon trop petit réduit la fiabilité (l. 10‑11).
\item \textbf{True.} Les sondages influencent le débat public avant les élections en mettant en lumière certains thèmes (l. 13‑15).
\item \textbf{False.} Le texte prévient que tous les sondages publiés ne sont pas fiables ; il faut vérifier la méthodologie (l. 17‑19).
\end{enumerate}
\end{corrige}

\begin{corrige}[Exercice 3 — Vocabulaire : mots croisés du thème]
\begin{enumerate}
\item The \textbf{results} of the poll will be announced on television tonight.
\item Sixty per cent of the \textbf{majority} think the government should build more schools. 
\item Our class conducted a \textbf{survey} on students' eating habits.
\item Journalists often \textbf{publish} poll results before elections.
\item In my \textbf{opinion}, young people should vote at eighteen.
\end{enumerate}
\emph{Note :} \eng{majority} correspond à « la majorité » (ici 60 %). Les mots du lexique sont utilisés dans leur forme correcte.
\end{corrige}

\begin{corrige}[Exercice 4 — Conjugaison : complète au bon temps]
\begin{enumerate}
\item Every year, thousands of Cameroonians \textbf{take} part in surveys. (présent simple : habitude)
\item The results of the poll \textbf{were published} in yesterday's newspaper. (passif passé simple : action terminée hier)
\item So far, the institute \textbf{has interviewed} two thousand respondents. (present perfect : action commencée dans le passé et continue)
\item Next month, a new survey \textbf{will be carried out} in our region. (futur passif : action future programmée)
\item People usually \textbf{answer} polls more honestly when they remain anonymous. (présent simple : vérité générale)
\end{enumerate}
\end{corrige}

\begin{corrige}[Exercice 5 — Traduction]
\begin{enumerate}
\item \eng{Polls allow citizens to express their opinion.}
\item \eng{The institute interviewed a sample of one thousand people in Douala.}
\item \eng{The margin of error of this poll is three per cent.}
\item \eng{Many young people refuse to answer phone surveys.}
\item \eng{The results will be published before the elections.}
\end{enumerate}
\end{corrige}

\begin{corrige}[Exercice 6 — Transformation de phrases]
\begin{enumerate}
\item \eng{500 voters in Yaoundé were interviewed by the institute.}
\item \eng{The journalist asked me if I trusted the results of that poll.}
\item \eng{The results are not reliable because the sample is too small.}
\item \eng{People do not participate in polls willingly.}
\item \eng{When was the survey conducted?}
\end{enumerate}
\end{corrige}

\begin{corrige}[Exercice 7 — Texte à trous]
\begin{textebox}[A survey in Buea — Corrigé]
Last month, a research centre \textbf{conducted} a survey on young people's reading habits in Buea. A \textbf{sample} of 800 students was selected from five secondary schools. Each student received a \textbf{questionnaire} with twenty questions. Most \textbf{respondents} answered honestly, so the results are considered \textbf{accurate}. The \textbf{margin} of error was only two per cent. The survey showed that the \textbf{majority} of students prefer reading on their phones, and in their \textbf{opinion}, libraries should offer more digital books. The organisers hope for a high \textbf{turnout} when the survey is repeated next year.
\end{textebox}
\end{corrige}

\begin{corrige}[Exercice 8 — Appariements : mots et définitions]
\begin{center}
\begin{tabular}{|c|c|}
\hline
\textbf{Mot} & \textbf{Définition} \\
\hline
1. respondent & b) a person who answers the questions of a survey \\
2. turnout & d) the number or percentage of people who take part in something \\
3. margin of error & g) the degree to which the results of a poll may be wrong \\
4. sample & e) a small group chosen to represent a larger population \\
5. census & c) an official count of all the people in a country \\
6. questionnaire & f) a written list of questions used in a survey \\
7. to conduct & h) to organise and carry out an activity \\
8. forecast & a) a prediction about what will happen, for example in an election \\
\hline
\end{tabular}
\end{center}
\end{corrige}

\begin{corrige}[Exercice 9 — Type Bac : compréhension écrite]
\textbf{Barème indicatif : 10 points au total (environ 1,5 pt par question, 1 pt pour les items lexicaux).}

\begin{enumerate}
\item \textbf{Where do interviewers in Lagos ask people about their voting intentions?} \\
\eng{At bus stops and in markets.} (l. 3‑4)
\item \textbf{Why does Mrs Adebayo refuse to answer pollsters' questions?} \\
\eng{Because the previous polls predicted a winner who actually lost, so she does not trust them.} (l. 5‑7)
\item \textbf{What condition must be met for a sample to represent a whole nation? Answer in your own words.} \\
\eng{The sample must be carefully selected and the margin of error clearly stated.} (l. 9‑10)
\item \textbf{Mention two reasons why poll results can be wrong, according to the text.} \\
\eng{Voters may change their minds at the last minute, and some hide their real intentions from interviewers.} (l. 11‑12)
\item \textbf{Give two positive effects of polls mentioned in the last paragraph.} \\
\eng{Polls push candidates to address voters' concerns (jobs, security, food prices) and encourage citizens to discuss politics openly.} (l. 14‑16)
\item \textbf{Find in the text : (a) a synonym of “figures”; (b) the opposite of “hide”.} \\
(a) \eng{numbers} (l. 2) \quad (b) \eng{show} / \eng{reveal} (l. 12 : « hide » → opposite « show »).
\end{enumerate}

\textbf{Points de vigilance :}
\begin{itemize}
\item Pour la Q3, reformuler sans copier-coller la phrase du texte.
\item Pour la Q6, le synonyme doit être un mot du texte (\"numbers\"), l'opposé peut être un verbe de sens contraire (\eng{show}, \eng{reveal}).
\item Respecter la limite de mots si imposée (ici pas de limite stricte, mais réponses concises).
\end{itemize}
\end{corrige}

\begin{corrige}[Exercice 10 — Type Bac : traduction]
\textbf{Traduction proposée :}
« Les sondeurs répondent qu'un échantillon soigneusement sélectionné peut effectivement représenter une nation entière, à condition que la marge d'erreur soit clairement indiquée. Ils admettent également que certains électeurs changent d'avis à la dernière minute, ou cachent leurs véritables intentions aux enquêteurs. »

\textbf{Barème indicatif : 5 points}
\begin{itemize}
\item \eng{Pollsters} → « Les sondeurs » / « Les instituts de sondage » (1 pt)
\item \eng{carefully selected sample} → « échantillon soigneusement sélectionné » (1 pt)
\item \eng{provided the margin of error is clearly stated} → « à condition que la marge d'erreur soit clairement indiquée / mentionnée » (1,5 pt)
\item \eng{admit that some voters change their minds at the last minute} → « admettent que certains électeurs changent d'avis à la dernière minute » (1 pt)
\item \eng{or hide their real intentions from interviewers} → « ou cachent leurs véritables intentions aux enquêteurs » (0,5 pt)
\end{itemize}

\textbf{Points de vigilance :}
\begin{itemize}
\item \eng{provided that} = « à condition que » + subjonctif en français.
\item \eng{margin of error} = terme technique « marge d'erreur ».
\item \eng{at the last minute} = « à la dernière minute » (locution figée).
\item \eng{hide} → « cachent » (pas « dissimulent » qui est plus soutenu mais acceptable).
\item Respecter le temps (présent) et la voix active/passive selon l'original.
\end{itemize}
\end{corrige}

\begin{corrige}[Exercice 11 — Type Bac : expression écrite]
\textbf{Barème indicatif : 10 points}
\begin{itemize}
\item Respect du nombre de mots (60‑80) : 2 pts
\item Utilisation d'au moins 3 mots du lexique (survey, respondents, sample, results, opinion) : 2 pts
\item Cohérence, organisation (but + opinion personnelle) : 3 pts
\item Correction grammaticale et orthographe : 3 pts
\end{itemize}

\textbf{Proposition de rédaction modèle (≈ 70 mots) :}
\eng{Our school is conducting a survey on students' use of social media. The purpose is to collect reliable data on how much time we spend online and which platforms we prefer. A representative sample of 200 respondents will answer a questionnaire. In my opinion, this survey is very useful because the results will help teachers understand our digital habits and design better media‑education programmes.}

\textbf{Points de vigilance :}
\begin{itemize}
\item Inclure explicitement le but (\eng{purpose}) et l'opinion (\eng{in my opinion}).
\item Utiliser correctement les termes : \eng{survey}, \eng{sample}, \eng{respondents}, \eng{results}, \eng{opinion}.
\item Vérifier l'accord des temps (présent pour le but, futur pour les résultats attendus).
\item Compter les mots : viser 65‑75 mots pour être dans la fourchette.
\end{itemize}
\end{corrige}

\newpage

\coursec{Fiche bilan}

\begin{popfigure}
\begin{tikzpicture}[mindmap, grow cyclic, every node/.style={concept}, text=white,
  root concept/.append style={concept color=popblue, font=\bfseries\small},
  level 1/.append style={level distance=3.1cm, sibling angle=60, font=\scriptsize},
  level 2/.append style={level distance=1.9cm, font=\tiny}]
\node[root concept, align=center]{Lesson 43\\Reading:\\Polls \& Surveys}
  child[concept color=poporange]{ node[align=center]{Key vocabulary\\poll, survey, sample} }
  child[concept color=popgreen]{ node[align=center]{Skimming\\gist, main idea} }
  child[concept color=poppurple]{ node[align=center]{Scanning\\dates, figures, names} }
  child[concept color=poppink]{ node[align=center]{Guessing\\words from context} }
  child[concept color=popgold]{ node[align=center]{Comprehension\\questions} }
  child[concept color=popblue!70]{ node[align=center]{Summary \&\\opinion} };
\end{tikzpicture}
\legende{Carte mentale de la leçon : lire des textes sur les sondages et les enquêtes d'opinion.}
\end{popfigure}

\begin{bilan}

\textbf{1. Lexique clé du thème (the language of polls)}

\begin{center}
\begin{tabularx}{\linewidth}{|X|X|}\hline
\rowcolor{popblueL} \textbf{English} & \textbf{Français} \\ \hline
a poll / an opinion poll & un sondage (d'opinion) \\ \hline
a survey & une enquête, une étude \\ \hline
to take part in a poll & participer à un sondage \\ \hline
to conduct / to carry out a survey & mener une enquête \\ \hline
a respondent / an interviewee & une personne interrogée \\ \hline
a sample (of the population) & un échantillon (de la population) \\ \hline
a questionnaire & un questionnaire \\ \hline
to answer / to fill in a questionnaire & répondre à / remplir un questionnaire \\ \hline
public opinion & l'opinion publique \\ \hline
the majority / the minority & la majorité / la minorité \\ \hline
the results / the findings & les résultats \\ \hline
a percentage & un pourcentage \\ \hline
to be in favour of / to be against & être favorable à / être opposé à \\ \hline
undecided voters & les électeurs indécis \\ \hline
a margin of error & une marge d'erreur \\ \hline
reliable / unreliable data & des données fiables / non fiables \\ \hline
\end{tabularx}
\end{center}

\textbf{2. Points de langue à connaître par cœur}

\begin{center}
\begin{tabularx}{\linewidth}{|X|X|}\hline
\rowcolor{popblueL} \textbf{Règle} & \textbf{Exemple} \\ \hline
Exprimer un pourcentage : \cle{per cent} en deux mots, après le chiffre & \eng{Sixty per cent of the respondents agree.} \\ \hline
\cle{the majority of} + nom pluriel + verbe au pluriel & \eng{The majority of young people trust online polls.} \\ \hline
Verbes d'opinion : \cle{to believe, to think, to claim, to argue, to state} + that & \eng{The report states that polls influence voters.} \\ \hline
Passif pour rapporter des résultats : \cle{be + participe passé} & \eng{Two thousand people were interviewed in Yaoundé.} \\ \hline
Comparatifs pour commenter des chiffres : \cle{higher than, fewer than, twice as many} & \eng{Fewer women than men took part in the survey.} \\ \hline
\end{tabularx}
\end{center}

\textbf{3. Méthodes express de lecture}
\begin{itemize}
\item \cle{Skimming} : lire le titre, la première phrase de chaque paragraphe et la conclusion pour trouver l'\textbf{idée générale} (gist) en moins de deux minutes.
\item \cle{Scanning} : repérer rapidement une \textbf{information précise} (chiffre, date, nom, lieu) sans tout lire ; laisser l'œil glisser sur le texte.
\item \cle{Guessing} : deviner le sens d'un mot inconnu grâce au \textbf{contexte}, à la \textbf{famille de mots} (\eng{respond} $\rightarrow$ \eng{respondent}) et aux \textbf{faux amis} à éviter.
\item Répondre aux questions : reformuler avec ses propres mots, citer une courte phrase du texte comme preuve.
\item Résumer : une phrase par paragraphe, au présent, sans exemples ni détails ; donner son opinion avec \eng{In my opinion...}, \eng{I believe that...}, \eng{As far as I am concerned...}.
\end{itemize}

\textbf{4. Expressions utiles pour la production}
\begin{itemize}
\item \eng{According to the text, ...} — D'après le texte, ...
\item \eng{The writer suggests that ...} — L'auteur laisse entendre que ...
\item \eng{This poll shows that most citizens ...} — Ce sondage montre que la plupart des citoyens ...
\item \eng{In my view, polls are useful because ...} — À mon avis, les sondages sont utiles parce que ...
\end{itemize}
\end{bilan}

\begin{aretenir}[Checklist avant le Bac]
\begin{itemize}
\item[$\square$] Je connais le vocabulaire des sondages : \eng{poll, survey, sample, respondent, questionnaire, margin of error}.
\item[$\square$] Je sais faire du \cle{skimming} pour trouver l'idée générale d'un texte en deux minutes.
\item[$\square$] Je sais faire du \cle{scanning} pour repérer vite un chiffre, une date ou un nom.
\item[$\square$] Je sais deviner le sens d'un mot inconnu grâce au contexte et à sa famille de mots.
\item[$\square$] Je me méfie des faux amis : \eng{a survey} n'est pas « un survol », \eng{to assist} n'est pas « assister à ».
\item[$\square$] J'écris \eng{per cent} en deux mots et je place le pourcentage correctement dans la phrase.
\item[$\square$] J'accorde le verbe après \eng{the majority of} + nom pluriel.
\item[$\square$] Je sais rapporter des résultats au passif : \eng{1,000 people were questioned}.
\item[$\square$] Je justifie mes réponses de compréhension par une courte citation du texte.
\item[$\square$] Je sais résumer un texte en trois à cinq phrases, sans recopier de longs passages.
\item[$\square$] Je sais donner mon opinion avec \eng{In my opinion}, \eng{I believe that}, suivi d'un argument.
\item[$\square$] Je relis ma production : orthographe, accords, ponctuation, majuscule à \eng{I}.
\end{itemize}
\end{aretenir}

\findecours

\end{document}
